\documentclass[11pt, openany]{book}
\usepackage{cpstyle}

% one card per library template (docs/generated/library_reference.tex)
\newtcolorbox{libentry}[6]{enhanced, breakable, colback=white, colframe=cpgray!35, boxrule=0.5pt, arc=2pt,
  fonttitle=\sffamily\bfseries\small, colbacktitle=cplight, coltitle=cpblue, title={#1},
  left=6pt, right=6pt, top=3pt, bottom=3pt, before skip=5pt, after skip=5pt,
  before upper={{\sffamily\scriptsize\key{#2}\hspace{0.8em}\texttt{\color{cporange}#3}\hspace{0.8em}%
    \textcolor{cpteal}{\faThumbtack\ #4}\hspace{0.8em}\textcolor{cpgray}{\faStopwatch\ #5}\hspace{0.8em}%
    \textcolor{cpgray}{\faFileCode[regular]\ \texttt{#6}}}\par\smallskip}}

\hypersetup{pdftitle={Competitive Programming Setup: User Guide}, pdfauthor={Kshitish Madbhavi}}

\begin{document}
\frontmatter
\begin{titlepage}
\begin{tikzpicture}[remember picture, overlay]
  \fill[cpblue] (current page.north west) rectangle ([yshift=-9.5cm]current page.north east);
  \fill[cpteal] ([yshift=-9.5cm]current page.north west) rectangle ([yshift=-9.9cm]current page.north east);
  % a little graph as decoration
  \begin{scope}[shift={([xshift=-5.2cm, yshift=-4.6cm]current page.north east)}, every node/.style={circle, fill=white, inner sep=2.2pt}]
    \node (a) at (0, 1.6) {}; \node (b) at (1.4, 2.2) {}; \node (c) at (2.6, 1.2) {}; \node (d) at (1.2, 0.3) {};
    \node (e) at (3.4, 2.6) {}; \node (f) at (0.2, -0.6) {}; \node[fill=cporange] (g) at (2.4, -0.4) {};
    \draw[white, thick] (a)--(b)--(c)--(d)--(a) (b)--(e) (c)--(e) (d)--(f) (d)--(g) (c)--(g);
  \end{scope}
  \node[anchor=north west, text=white, font=\sffamily\bfseries\fontsize{34}{40}\selectfont, align=left]
    at ([xshift=2.2cm, yshift=-2.6cm]current page.north west) {Competitive\\Programming\\Setup};
  \node[anchor=north west, text=cpteal!40!white, font=\sffamily\Large, align=left]
    at ([xshift=2.2cm, yshift=-7.6cm]current page.north west) {User Guide \textperiodcentered{} VS Code + CPH};
  \node[anchor=north west, text=cpgray, font=\sffamily\large, align=left, text width=15cm]
    at ([xshift=2.2cm, yshift=-11.2cm]current page.north west)
    {Everything installed in \texttt{Desktop/competitive}: the template, the 175-template library and its
     keyboard shortcuts, the CPH integration, the test runner, the stress tester and the tools behind them,
     and how to use every piece of it.};
  \node[anchor=south west, text=cpblue, font=\sffamily\bfseries\large]
    at ([xshift=2.2cm, yshift=3.2cm]current page.south west) {Kshitish Madbhavi};
  \node[anchor=south west, text=cpgray, font=\sffamily]
    at ([xshift=2.2cm, yshift=2.5cm]current page.south west) {Companion volume: \emph{The Competitive Programming Handbook} (ALGORITHMS.pdf)};
\end{tikzpicture}
\end{titlepage}

\tableofcontents
\mainmatter
\chapter{Overview}\label{ch:overview}

This guide documents the competitive programming environment set up in the \file{Desktop/\allowbreak competitive} folder.
It is built around one idea: \textbf{the CPH extension stays in charge of test cases}. You fetch a problem, CPH creates
the file from your template and stores the samples, you write code, and you press CPH's \emph{Run All} button.
Everything else in this setup either makes that loop faster or explains what CPH cannot.

\section{The big picture}

\begin{figure}[H]
\centering
\begin{tikzpicture}[
  box/.style={draw=cpblue, fill=cpblue!6, rounded corners=3pt, thick, minimum width=3.2cm, minimum height=1.25cm,
    align=center, font=\sffamily\small},
  tool/.style={draw=cporange, fill=cporange!8, rounded corners=3pt, thick, minimum width=3.0cm, minimum height=1.05cm,
    align=center, font=\sffamily\small},
  arr/.style={->, thick, cpgray},
  lab/.style={font=\sffamily\scriptsize, text=cpgray, align=center}]
  \node[box] (browser) at (0, 0) {Problem page\\\scriptsize Codeforces, USACO, CSES};
  \node[box] (cc) at (4.5, 0) {Competitive\\Companion};
  \node[box, fill=cpteal!10, draw=cpteal] (cph) at (9, 0) {\textbf{CPH}\\\scriptsize stores the samples};
  \node[box] (file) at (9, -2.7) {\file{A\_Name.cpp}\\\scriptsize from \file{template.cpp}};
  \node[box] (editor) at (4.5, -2.7) {You write code\\\scriptsize templates, \texttt{dbg}};
  \node[tool] (f5) at (0, -2.7) {\key{F5}\ CPH tests\\\scriptsize crash reason, memory};
  \node[box, fill=cpgreen!10, draw=cpgreen] (run) at (4.5, -5.4) {CPH \emph{Run All}\\\scriptsize \key{Ctrl+Alt+B}};
  \node[box] (submit) at (9, -5.4) {Submit\\\scriptsize the file as it is};
  \node[tool] (stress) at (0, -5.4) {Stress test\\\scriptsize \key{Ctrl+Alt+X}};
  \draw[arr] (browser) -- node[above, lab] {click} (cc);
  \draw[arr] (cc) -- node[above, lab] {sends problem} (cph);
  \draw[arr] (cph) -- node[right, lab] {creates the file,\\fills the header\\and USACO I/O} (file);
  \draw[arr] (file) -- node[above, lab] {opens} (editor);
  \draw[arr] (editor) -- node[right, lab] {save} (run);
  \draw[arr] (run) -- node[above, lab] {all pass} (submit);
  \draw[arr, dashed] (editor) -- node[above=1pt, lab] {verdict is\\a mystery} (f5);
  \draw[arr, dashed] (f5) -- node[left, lab] {still stuck} (stress);
  \draw[arr, dashed] (stress) -- node[below=1pt, lab] {paste the\\failing input} (run);
\end{tikzpicture}
\caption{The daily loop. Solid arrows are the normal path; dashed arrows are what you do when a verdict surprises you.}
\end{figure}

\section{Quick start in 60 seconds}

\begin{enumerate}[label=\textbf{\color{cpteal}\arabic*.}]
  \item Open the \file{competitive} folder in VS Code. It uses the \emph{competitive} profile automatically.
  \item On a problem page, click the green Competitive Companion button. CPH creates \file{A\_Problem\_Name.cpp}
        from the template, fills in the problem name, link and time, and (for USACO) switches on file I/O.
  \item Write your solution in \cc{solve()}. Type \texttt{cp} to see every library snippet, or press
        \key{Ctrl+Alt+T} followed by a key (for example \key{U} for a DSU) to paste one at the cursor.
  \item Use \cc{dbg(x, v)} freely: it prints to CPH's \emph{Standard Error} box and vanishes on the judge.
  \item Press \key{Ctrl+Alt+B} (or CPH's \emph{Run All} button). If a verdict is confusing, press \key{F5}: the same tests
        run in the terminal with the crash reason, time and peak memory of each one.
  \item When all tests pass, submit the file exactly as it is. Nothing needs to be deleted or uncommented.
\end{enumerate}

\begin{tip}
The one-page summary of every shortcut is \file{CHEATSHEET.md} in the folder root. This guide is the long version;
the algorithms behind the templates are taught in the companion handbook, \emph{ALGORITHMS.pdf}.
\end{tip}

\section{What is where}

\newcommand{\desc}[1]{{\normalfont\sffamily\scriptsize\color{cpgray}--- #1}}
\begin{figure}[H]
\centering
\begin{forest}
  for tree={font=\ttfamily\small, grow'=0, child anchor=west, parent anchor=south, anchor=west, calign=first,
    edge path={\noexpand\path [draw, cpgray, \forestoption{edge}] (!u.south west) +(7.5pt,0) |- (.child anchor)\forestoption{edge label};},
    fit=band, s sep=1pt, before computing xy={l=15pt}}
  [{competitive/}
    [{template/template.cpp\quad\desc{the file CPH copies for every new problem}}]
    [{template/template.cpp.bak\quad\desc{your original template}}]
    [{include/debug.h\quad\desc{the dbg() printer (local builds only)}}]
    [{include/bits/\quad\desc{precompiled bits/stdc++.h (generated)}}]
    [{library/\quad\desc{129 templates in 10 category folders}}]
    [{tools/cp.py\quad\desc{test runner, stress tester, quickstart, setup, docs}}]
    [{stress/gen.cpp, brute.cpp\quad\desc{inputs for the stress tester}}]
    [{.vscode/tasks.json, settings.json\quad\desc{tasks behind F5 etc., IntelliSense}}]
    [{.cph/\quad\desc{CPH's test cases, one .prob file per problem (hidden)}}]
    [{docs/\quad\desc{this guide and the algorithms handbook (LaTeX + PDF)}}]
    [{archive/YYYY-MM/\quad\desc{old solutions, created by the archive shortcut}}]
    [{CHEATSHEET.md\quad\desc{one-page shortcut table}}]
  ]
\end{forest}
\caption{The folder layout. Build output (\texttt{.build/}, \texttt{*.exe}, \texttt{*.bin}) and \texttt{.cph/} are hidden from VS Code's explorer.}
\end{figure}

\section{What was changed outside the folder}

Keyboard shortcuts and snippets can only live in VS Code's user settings, so a few files outside the folder were
changed. All of them belong to the \emph{competitive} VS Code profile only; your default profile is untouched.

\begin{center}
\begin{tabularx}{\linewidth}{>{\raggedright\arraybackslash}p{5.2cm}L}
\toprule
\thd{File (in the profile folder)} & \thd{What changed} \\
\midrule
\file{keybindings.json} & One generated block, between \texttt{>>> cp-library shortcuts} and \texttt{<<<} markers, with the
  Ctrl+Alt+T / Ctrl+Alt+Y / Ctrl+Alt+U shortcuts, Ctrl+Alt+X and Ctrl+F5. Your own bindings above it are untouched.
  Backup: \file{keybindings.json.bak-cp-setup}. \\
\file{snippets/\allowbreak cp-library.code-snippets} & New file: the 175 snippets. \\
\file{settings.json} & Four CPH settings: compiler flags, template path, template variables on, stderr ignored for
  verdicts. Backup: \file{settings.json.bak-cp-setup}. \\
\midrule
TinyTeX & LaTeX packages were added (fonts, listings, forest, \ldots) to compile these PDFs. Additive only. \\
\bottomrule
\end{tabularx}
\end{center}

The profile folder is \file{\%APPDATA\%/Code/User/profiles/5d3986b6}. Re-running \emph{CP: Setup} rewrites only the
generated block and the four CPH settings, so it is always safe to run again.

\section{The four shortcuts that were broken before}

Your profile already had \key{F5}, \key{Ctrl+Shift+N}, \key{Ctrl+Shift+U} and \key{Ctrl+Shift+Alt+C} bound to tasks
named \texttt{cf-run}, \texttt{CP: Quickstart}, \texttt{CP: Quickstart USACO} and \texttt{CP: Archive Old Files},
but no such tasks existed anywhere, so the keys did nothing. They are now defined in \file{.vscode/tasks.json} and
work as described in Chapter~\ref{ch:testing} and Section~\ref{sec:quickstart}.

\section{How this setup is tied to CPH}\label{sec:cph-tie}

CPH is the centre of the workflow, so every piece was checked against CPH's own source code:

\begin{itemize}
  \item \textbf{CPH compiles with your flags.} CPH runs \texttt{g++ file.cpp -o file.bin \emph{Args} -D DEBUG -D CPH}.
        The \emph{Args} setting is written by \emph{CP: Setup} from one list in \file{tools/cp.py}, so the precompiled
        header, the 256\,MB stack, static linking and \cc{-DLOCAL} all apply to CPH runs.
        (The extra \texttt{-D DEBUG -D CPH} were verified not to disturb the precompiled header.)
  \item \textbf{CPH fills the template.} With template variables enabled, CPH replaces \texttt{\$name\$},
        \texttt{\$group\$}, \texttt{\$url\$}, \texttt{\$date\$}, \texttt{\$time\$}, \texttt{\$input\$} and
        \texttt{\$output\$} when it creates a file. That is how the header and the USACO file I/O are filled in.
  \item \textbf{CPH shows \cc{dbg} output.} CPH treats any stderr output as a failed test unless
        \texttt{ignoreSTDERROR} is on. It is now on, so \cc{dbg} output appears in the \emph{Standard Error}
        box while verdicts depend only on stdout and the exit code.
  \item \textbf{F5 reads CPH's tests.} The terminal runner reads the same \file{.cph/*.prob} file CPH uses.
  \item \textbf{The stress tester checks your brute force against CPH's samples} before it trusts it, and puts a
        failing input on the clipboard ready to paste into CPH as a new test case.
\end{itemize}

\begin{fact}[Honest limits]
Nothing is lost when you use only the CPH button. What CPH cannot do by design: tell you \emph{why} a program
crashed (it just fails the test), show memory use, or accept test cases from outside tools while its panel is open.
F5 covers the first two; the clipboard hand-off covers the third. The checked STL (\texttt{-D\_GLIBCXX\_DEBUG}) makes
local running times slower than the judge's, so treat CPH's times as upper bounds.
\end{fact}

\chapter{The Daily Workflow}\label{ch:workflow}

\section{Solving a problem, step by step}

\begin{figure}[H]
\centering
\begin{tikzpicture}[
  stp/.style={draw=cpteal, fill=cpteal!7, rounded corners=4pt, thick, text width=3.35cm, minimum height=2.15cm,
    align=center, font=\sffamily\small},
  num/.style={circle, fill=cpteal, text=white, font=\sffamily\bfseries\small, inner sep=1.5pt, minimum size=5.5mm},
  arr/.style={->, very thick, cpteal!70}]
  \foreach \px/\ti/\tx [count=\k] in {
      0/Fetch/{Competitive Companion button on the problem page},
      4.1/Read/{problem text, constraints, sample I/O},
      8.2/Code/{\texttt{solve()}; paste templates with \key{Ctrl+Alt+T}},
      12.3/Test/{CPH \emph{Run All} \key{Ctrl+Alt+B}}} {
    \node[stp] (s\k) at (\px, 0) {\textbf{\ti}\\[2pt]\scriptsize \tx};
    \node[num] at (s\k.north west) {\k};
  }
  \foreach \px/\ti/\tx [count=\k from 5] in {
      12.3/Explain/{\key{F5}: crash reason, time, memory per test},
      8.2/Hunt/{\key{Ctrl+Alt+X}: stress test vs a brute force},
      4.1/Fix/{add the failing input to CPH, fix, rerun},
      0/Submit/{the file exactly as it is}} {
    \node[stp] (s\k) at (\px, -3.1) {\textbf{\ti}\\[2pt]\scriptsize \tx};
    \node[num] at (s\k.north west) {\k};
  }
  \draw[arr] (s1) -- (s2); \draw[arr] (s2) -- (s3); \draw[arr] (s3) -- (s4);
  \draw[arr] (s4) -- node[right, font=\sffamily\scriptsize, text=cpgray] {fails} (s5);
  \draw[arr] (s5) -- (s6); \draw[arr] (s6) -- (s7);
  \draw[arr] (s7) -- node[above, font=\sffamily\scriptsize, text=cpgray] {passes} (s8);
  \draw[arr, cpgreen] (s4.north) to[out=90, in=90, looseness=0.35] node[above, font=\sffamily\scriptsize, text=cpgreen] {all tests pass: go straight to submit} (s1.north);
\end{tikzpicture}
\caption{Steps 5--7 are only needed when a test fails and the reason is not obvious.}
\end{figure}

\subsection{1. Fetch}
Click the Competitive Companion button in the browser. CPH receives the problem (name, contest, link, samples, limits
and whether it reads from a file) and creates \file{A\_Problem\_Name.cpp} from \file{template/template.cpp}.
The file opens with its header already filled in:

\begin{code}
/**
 *    author:  Kshitish Madbhavi
 *    problem: A. Graph
 *    source:  Codeforces - Baltic Olympiad in Informatics 2020, Day 2 (IOI, Unofficial Mirror Contest, Unrated)
 *    url:     https://codeforces.com/contest/1387/problem/A
 *    created: 2026-09-29 20:31:55
 */
\end{code}

For a USACO problem the same step also switches on file input and output; you do not touch anything
(Section~\ref{sec:usaco}).

\subsection{2. Read}
Look at the constraints before thinking about an algorithm: they tell you the complexity you can afford
(ALGORITHMS.pdf, Chapter~2 has the full table). Decide whether the input has several test cases. If it starts
with ``the first line contains $t$'', uncomment \cc{cin >> t;} in \cc{main} (put the cursor on that line and press
\key{Ctrl+/}).

\subsection{3. Code}
Write the solution in \cc{solve()}. Three ways to pull in library code, all of which paste at the cursor:

\begin{center}
\begin{tabularx}{\linewidth}{p{4.3cm}L}
\toprule
\thd{Method} & \thd{Example} \\
\midrule
Chord shortcut & \key{Ctrl+Alt+T} then \key{J} pastes Dijkstra; \key{Ctrl+Alt+Y} then \key{H} pastes HLD. \\
Prefix + \key{Tab} & type \texttt{cpdsu} and press \key{Tab}. Every prefix starts with \texttt{cp}, so typing \texttt{cp} lists all of them. \\
Picker & \key{Ctrl+Alt+T} then \key{Space}: a searchable list of every template. \\
\bottomrule
\end{tabularx}
\end{center}

Chapter~\ref{ch:library} explains the library and has keyboard maps of every chord.

\subsection{4. Test}
Press \key{Ctrl+Alt+B} or CPH's \emph{Run All} button. CPH compiles with the flags from Chapter~\ref{ch:compiler}
(about 0.4--0.9\,s thanks to the precompiled header) and runs every sample. \cc{dbg} output appears in each test's
\emph{Standard Error} box without affecting the verdict. Add your own tests with CPH's \emph{New Testcase} button;
edge cases (smallest $n$, largest $n$, all equal values, negative numbers) catch most wrong answers.

\subsection{5--7. When a test fails}
\begin{itemize}
  \item \textbf{Wrong answer and you can see why}: fix it and rerun.
  \item \textbf{Runtime error, or a time you do not understand}: press \key{F5}. The same CPH tests run in the terminal
        and each failure shows the input, expected output, your output, your stderr and a plain-language crash
        reason such as ``stack overflow'' (Section~\ref{sec:f5}).
  \item \textbf{The samples pass but the judge says wrong answer}: write a slow brute force in \file{stress/brute.cpp}
        and a random input generator in \file{stress/gen.cpp}, then press \key{Ctrl+Alt+X}. The stress tester finds a
        small input where your answer differs and copies it to the clipboard; paste it into CPH as a new test
        (Section~\ref{sec:stress}).
\end{itemize}

\subsection{8. Submit}
Submit the file unchanged: \cc{dbg} lines compile to nothing on the judge, and USACO file I/O is already on. CPH's
\emph{Submit} (\key{Ctrl+Alt+S}) sends Codeforces submissions as \emph{GNU G++23 14.2 (64 bit)}. Your code is still
strictly C++17: that is what CPH and IntelliSense compile it as, and every library template is checked with
\texttt{-std=c++17 -Wpedantic}. G++23 was kept because it is the newest 64-bit compiler Codeforces offers and
accepts C++17 code unchanged.

\section{Keyboard shortcuts}

\begin{center}
\begin{tabularx}{\linewidth}{>{\raggedright\arraybackslash}p{4.4cm}L}
\toprule
\thd{Keys} & \thd{Action} \\
\midrule
\key{Ctrl+Alt+B} & CPH: compile and run all test cases (CPH's own shortcut). \\
\key{Ctrl+Alt+S} & CPH: submit to Codeforces. \\
\key{Ctrl+Alt+T}, \emph{key} & paste a core template (Bronze to Gold level). \\
\key{Ctrl+Alt+Y}, \emph{key} & paste an advanced template (Platinum and beyond). \\
\key{Ctrl+Alt+U}, \emph{key} & paste a specialist template (polynomials, matchings, cuts, exotic structures). \\
\key{Ctrl+Alt+T}, \key{Space} & searchable template picker. \\
\key{F5} & run the CPH test cases in the terminal with diagnostics. \\
\key{Ctrl+F5} & run with input you type in the terminal. \\
\key{Ctrl+Alt+X} & stress test the current file. \\
\key{Ctrl+Shift+N} & new problem file from the template (asks for a name). \\
\key{Ctrl+Shift+U} & new USACO problem (asks for the file name, e.g.\ \texttt{moocast}). \\
\key{Ctrl+Shift+Alt+C} & archive solutions untouched for a day into \file{archive/YYYY-MM/}. \\
\key{Ctrl+Alt+N} & Code Runner: same as \key{Ctrl+F5}. \\
\bottomrule
\end{tabularx}
\end{center}

\section{Multi-test problems}

Competitive Companion does not tell CPH whether a problem has several test cases (it always reports
``single''), so this cannot be detected automatically. The template therefore reads one test by default:

\begin{code}
int t = 1;
// cin >> t;
while (t--) solve();
\end{code}

For a multi-test problem, uncomment the middle line. Remember that everything \cc{solve()} uses must be reset
between test cases: declare arrays inside \cc{solve()} (sized by that test's $n$) rather than globally.

\begin{pitfall}[Resetting global state]
The most common multi-test bug is a global array or counter that keeps values from the previous test.
Either declare it inside \cc{solve()} or clear exactly the part you used; never \cc{memset} a $10^6$-element
array $10^5$ times, which is $10^{11}$ operations.
\end{pitfall}

\section{USACO problems}\label{sec:usaco}

USACO reads from \texttt{name.in} and writes \texttt{name.out}. When Competitive Companion sends a USACO problem it
tells CPH the file names, and CPH writes them into the template (Section~\ref{sec:fileio} shows exactly how).
On the USACO judge your program opens the files; locally, CPH and F5 keep feeding it standard input, so the
samples still run in CPH. There is nothing to uncomment.

For a USACO problem typed in by hand, press \key{Ctrl+Shift+U} and enter the file name (\texttt{moocast}): the new file
behaves exactly like one created by CPH.

\section{Quickstart and archiving}\label{sec:quickstart}

\begin{itemize}
  \item \key{Ctrl+Shift+N} asks for a name, turns it into a safe file name (\texttt{B. Two Sum} becomes
        \file{B\_Two\_Sum.cpp}), fills the template exactly as CPH would and opens the file. Use it for problems
        without Competitive Companion support. If the file exists it is simply opened.
  \item \key{Ctrl+Shift+U} does the same for USACO and fills in \texttt{name.in} / \texttt{name.out}.
  \item \key{Ctrl+Shift+Alt+C} moves every solution in the folder root that has not been modified for a day into
        \file{archive/YYYY-MM/} (by the month it was last edited), deletes stray \texttt{.exe} / \texttt{.bin} files,
        and moves each problem's CPH test cases along with it. CPH names a problem's test file after the MD5 hash of
        the source path, so the test file is renamed for the new path and CPH still finds your tests if you open the
        archived file. Nothing is ever deleted except build output.
\end{itemize}

\chapter{The Template, Line by Line}\label{ch:template}

\file{template/template.cpp} is the file CPH copies for every new problem. It is strictly C++17, compiles without a
single warning under \texttt{-Wall -Wextra -Wshadow -Wpedantic}, and was tested three ways: as CPH builds it locally,
as a C++17 judge builds it, and as a C++20 judge builds it. This chapter walks through it top to bottom.
Your original template is kept as \file{template/template.cpp.bak}.

\section{Header}

\begin{code}
/**
 *    author:  Kshitish Madbhavi
 *    problem: $name$
 *    source:  $group$
 *    url:     $url$
 *    created: $date$ $time$
 */
\end{code}

The \texttt{\$...\$} words are CPH template variables. When CPH creates a file it replaces each one with the
corresponding field of the problem it received, so every solution records where it came from. Files made with
\key{Ctrl+Shift+N} are filled the same way by \file{tools/cp.py}; the \texttt{cptemplate} snippet fills them with
VS Code's date and file-name variables.

\section{Includes and the debug hook}

\begin{code}
#include <bits/stdc++.h>
using namespace std;

#ifdef LOCAL
#include "debug.h"
#else
#define dbg(...)
#endif
\end{code}

\cc{LOCAL} is defined by every local build (CPH, F5, the stress tester) and never by a judge. Locally,
\file{include/debug.h} provides \cc{dbg(...)}; on the judge \cc{dbg(...)} is defined as nothing, so
\cc{dbg(a, b);} becomes an empty statement. Chapter~\ref{ch:testing} shows what \cc{dbg} prints.

\section{Types}

\begin{code}
#define int long long

using ll = long long;
using ld = long double;
using pii = pair<int, int>;
using pll = pair<ll, ll>;
using vi = vector<int>;
using vvi = vector<vi>;
template <class T> using minpq = priority_queue<T, vector<T>, greater<T>>;
\end{code}

\cc{#define int long long} turns every \cc{int} in your code into a 64-bit integer, which removes the most common
source of wrong answers: overflow of intermediate values such as \cc{a * b} when both are up to $10^9$.

\begin{fact}[Fixed from the old template]
In the old template \cc{#define int long long} came \emph{after} the aliases, so \cc{pii} was still
\cc{pair<int,int>} with 32-bit ints. Pushing a \cc{long long} distance into a \cc{minpqpii} silently truncated it,
a classic source of mysterious Dijkstra failures. The define now comes first, so every alias is 64-bit.
\end{fact}

Consequences of \cc{#define int long long} to remember:
\begin{itemize}
  \item \cc{main} must be declared as \cc{signed main()} (it is), because \cc{long long main()} is illegal.
  \item \cc|printf("%d", x)| is wrong for an \cc{int}; use \cc{cout} or \cc|%lld|.
  \item \cc{max(x, 0)} fails to compile when \cc{x} is \cc{int}, because \cc{0} is a real 32-bit \cc{int}
        literal. Write \cc{max<int>(x, 0)}, \cc{max(x, 0LL)} or use \cc{chmax} below (it accepts mixed types).
  \item \cc{__builtin_popcount(x)} truncates to 32 bits; use \cc{__builtin_popcountll}.
  \item Memory doubles for \cc{vector<int>}. $10^7$ \cc{int}s are now 80\,MB instead of 40\,MB.
\end{itemize}

\cc{minpq<T>} is a min-heap: \cc|minpq<pii> pq; pq.push({dist, v});| pops the smallest pair first. It replaces the
four separate \cc{minpqi}/\cc{minpql}/\cc{minpqpii}/\cc{minpqpll} aliases of the old template.

\section{Macros}

\begin{code}
#define all(x) (x).begin(), (x).end()
#define rall(x) (x).rbegin(), (x).rend()
#define sz(x) (int)(x).size()
#define rep(...)   // rep(i, n): 0..n-1   rep(i, a, b): a..b-1
#define per(...)   // same ranges, backwards
\end{code}

\begin{center}
\begin{tabularx}{\linewidth}{>{\ttfamily}p{4.2cm}>{\ttfamily}p{5.8cm}L}
\toprule
\thd{You write} & \thd{It means} & \thd{Typical use} \\
\midrule
sort(all(v)) & sort(v.begin(), v.end()) & any STL algorithm on a whole container \\
sort(rall(v)) & sort(v.rbegin(), v.rend()) & sort in decreasing order \\
sz(v) & (int)v.size() & avoids signed/unsigned warnings and \texttt{v.size() - 1} underflow \\
rep(i, n) & for (int i = 0; i < n; i++) & the everyday loop \\
rep(i, 1, n + 1) & for (int i = 1; i < n + 1; i++) & 1-indexed loops \\
per(i, n) & for (int i = n - 1; i >= 0; i--) & reverse loops \\
per(i, a, b) & for (int i = b - 1; i >= a; i--) & reverse over $[a, b)$ \\
\bottomrule
\end{tabularx}
\end{center}

\cc{rep} picks its two- or three-argument form by counting arguments (the \cc{PICK3} helper macro), a standard
C++17-compatible preprocessor trick.

\begin{pitfall}[Why \texttt{sz} matters]
\cc{v.size()} is unsigned. If \cc{v} is empty, \cc{v.size() - 1} is $2^{64}-1$, so
\cc{for (int i = 0; i < v.size() - 1; i++)} runs forever. \cc{sz(v) - 1} is $-1$ and the loop is skipped.
Library code avoids the \cc{sz} name for anything else, because the macro would swallow it.
\end{pitfall}

\section{Constants}

\begin{code}
constexpr ll INF = 4e18;
constexpr int MOD = 1'000'000'007;  // 998'244'353
\end{code}

\cc{INF} is large enough to mean ``unreachable'' yet \cc{INF + INF} still fits in a \cc{long long}
($8\times10^{18} < 9.22\times10^{18}$). \cc{MOD} is used by every modular template in the library (\texttt{cpmint},
\texttt{cpncr}, \ldots); switch the comment when a problem wants 998\,244\,353.

\section{Helpers}

\begin{code}
chmin(a, b);  chmax(a, b);   // a = min(a, b) / max(a, b); true if a changed; mixed types OK
cdiv(a, b);   fdiv(a, b);    // ceil(a / b) and floor(a / b), correct for negative numbers too
cin >> v;                    // reads a whole vector (sized beforehand) or a pair
cout << v;                   // prints a vector separated by spaces
print(a, b, c);              // prints "a b c" and a newline
yn(ok);                      // prints YES or NO
mex(v);                      // smallest non-negative integer not in v, O(n)
Mex m(n); m.add(x); m.remove(x); m.get();   // mex of a changing multiset (at most n values), O(log n)
\end{code}

\begin{center}
\begin{tabular}{>{\ttfamily}l>{\ttfamily}r>{\ttfamily}r>{\ttfamily}r}
\toprule
\thd{expression} & \thd{C++ \texttt{/}} & \thd{cdiv} & \thd{fdiv} \\
\midrule
7 / 2 & 3 & 4 & 3 \\
-7 / 2 & -3 & -3 & -4 \\
7 / -2 & -3 & -3 & -4 \\
-8 / -3 & 2 & 3 & 2 \\
\bottomrule
\end{tabular}
\end{center}

C++ division truncates toward zero, which is neither floor nor ceiling for negative results; \cc{cdiv} and
\cc{fdiv} were checked on all sign combinations. \cc{chmin}/\cc{chmax} take two type parameters so
\cc{chmax(ans, 0)} compiles even though \cc{0} is a plain \cc{int}.

\section{Automatic USACO file I/O}\label{sec:fileio}

\begin{code}
void file_io() {
#ifndef LOCAL
    auto file = [](const string& spec) { /* extract "fileName":"..." */ };
    string in = file(R"($input$)"), out = file(R"($output$)");
    if (!in.empty()) freopen(in.c_str(), "r", stdin);
    if (!out.empty()) freopen(out.c_str(), "w", stdout);
#endif
}
\end{code}

\begin{figure}[H]
\centering
\begin{tikzpicture}[
  bx/.style={draw=cpblue, fill=cpblue!5, rounded corners=3pt, thick, align=left, font=\ttfamily\scriptsize, inner sep=5pt},
  arr/.style={->, very thick, cpteal}, lab/.style={font=\sffamily\scriptsize, text=cpgray, align=center}]
  \node[bx] (json) at (0, 0) {"input": \{"type": "file",\\ \ \ \ \ \ \ \ \ \ "fileName": "moocast.in"\}\\
    "output": \{"type": "file",\\ \ \ \ \ \ \ \ \ \ \ "fileName": "moocast.out"\}};
  \node[bx, draw=cpteal, fill=cpteal!5] (tmpl) at (0, -2.6) {string in = file(R"(\$input\$)");};
  \node[bx, draw=cporange, fill=cporange!6] (res) at (8.1, -2.6) {string in = file(R"("type":"file",\\
    \ \ \ \ \ \ \ \ \ \ \ \ \ \ \ \ \ \ "fileName":"moocast.in")");};
  \node[bx, draw=cpgreen, fill=cpgreen!6] (judge) at (8.1, 0) {judge: freopen("moocast.in") ...\\locally (-DLOCAL): skipped, CPH uses stdin};
  \node[lab, above=2pt of json] {Competitive Companion sends};
  \node[lab, below=2pt of tmpl] {\file{template.cpp}};
  \draw[arr] (json) -- node[left, lab] {CPH} (tmpl);
  \draw[arr] (tmpl) -- node[above, lab] {\$input\$ replaced} (res);
  \draw[arr] (res) -- node[right, lab] {at run time} (judge);
\end{tikzpicture}
\caption{How a USACO problem switches on file I/O without any manual step.}
\end{figure}

CPH replaces \texttt{\$input\$} with the problem's input description, minus its outer braces. For USACO that is
\texttt{"type":"file","fileName":"moocast.in"}, so \cc{file()} finds a file name and \cc{freopen} runs; for every
other judge it is \texttt{"type":"stdin"}, no file name is found and nothing happens. The raw string literal
\cc{R"(...)"} means the quotes need no escaping. The whole function is compiled only when \cc{LOCAL} is
\emph{not} defined, so CPH keeps running your samples through standard input.

\begin{tip}
If a file was created some other way and you need file I/O, change the placeholder by hand to
\texttt{R"("fileName":"name.in")"} (and the output one to \texttt{name.out}), or recreate the file with
\key{Ctrl+Shift+U}.
\end{tip}

\section{solve() and main()}

\begin{code}
void solve() {

}

signed main() {
    ios::sync_with_stdio(false);
    cin.tie(nullptr);
    cout << fixed << setprecision(10);
    file_io();

    int t = 1;
    // cin >> t;
    while (t--) solve();
}
\end{code}

\begin{itemize}
  \item \cc{sync_with_stdio(false)} and \cc{cin.tie(nullptr)} make \cc{cin}/\cc{cout} as fast as \cc{scanf}. After them,
        do not mix \cc{cin} with \cc{scanf} or \cc{cout} with \cc{printf}.
  \item \cc{fixed << setprecision(10)} prints floating-point answers with ten decimals, which satisfies every
        ``absolute or relative error $10^{-6}$'' checker. Integers are unaffected.
  \item Print \cc|'\n'|, not \cc{endl}: \cc{endl} flushes the output every time and can turn a fast program into
        a TLE. The one exception is interactive problems, which need the flush (\texttt{cpinteractive}).
\end{itemize}

\chapter{Testing and Debugging}\label{ch:testing}

\section{Debug printing with dbg()}\label{sec:dbg}

\cc{dbg(...)} prints any number of expressions with their names and the line number, to standard error:

\begin{code}
int x = 3; vector<int> v{1, 2};
map<string, pair<int, char>> mp{{"k", {7, 'z'}}};
dbg(x, v, mp);
dbg(v[0] + v[1]);
dbg();                                  // just "reached": is this line executed?
\end{code}
\begin{shell}
[L42] x = 3 | v = {1, 2} | mp = {("k", (7, 'z'))}
[L43] v[0] + v[1] = 3
[L44] reached
\end{shell}

\begin{center}
\begin{tabularx}{\linewidth}{p{4.8cm}L}
\toprule
\thd{Type} & \thd{Printed as} \\
\midrule
numbers, \texttt{\_\_int128} & \texttt{42}, \texttt{-1267650600228229401496703205376} \\
\texttt{bool}, \texttt{char}, strings & \texttt{true}, \texttt{'z'}, \texttt{"text"} \\
pairs and tuples & \texttt{(1, "s", 2.5)} \\
any container, nested & \texttt{\{\{1, 2\}, \{3\}\}}; maps as \texttt{\{(key, value), ...\}} \\
\texttt{vector<bool>} & \texttt{\{true, false\}} \\
\texttt{stack}, \texttt{queue}, \texttt{priority\_queue} & contents from the top/front, the original is untouched \\
\texttt{bitset} & \texttt{0101} \\
\bottomrule
\end{tabularx}
\end{center}

Where the output appears: in CPH it is the \emph{Standard Error} box under each test; with \key{F5} it is printed
after the test's output when that test fails; with \key{Ctrl+F5} it goes straight to the terminal. The line is
flushed immediately, so the last \cc{dbg} before a crash is always visible, which makes \cc{dbg()} a quick way
to find the line that crashes.

\begin{fact}[Why dbg never fails a test in CPH]
CPH marks a test as failed whenever the program writes anything to standard error, unless the
\texttt{cph.general.ignoreSTDERROR} setting is on. \emph{CP: Setup} turns it on. Verdicts then depend on the output and
the exit code only, and crashes are still caught through the exit code.
\end{fact}

\begin{tip}
\cc{dbg} is compiled out on the judge, so you never need to delete debug lines before submitting. Heavy \cc{dbg}
calls inside a hot loop still slow down local runs, so comment those out when timing.
\end{tip}

\section{F5: CPH's tests with diagnostics}\label{sec:f5}

\key{F5} (task \texttt{cf-run}) runs \file{tools/cp.py test} on the current file. It reads the same test cases CPH
shows, from \file{.cph/.\textless file\textgreater\_\textless hash\textgreater.prob}, compiles with CPH's flags and
runs every test with a timeout:

\begin{shell}
[cp] compiled tricky.cpp in 0.4 s
[cp] tricky.cpp: 3 CPH test case(s), time limit 1000 ms
[cp] test 1: PASS  153 ms, 0.9 MB
[cp] test 2: WRONG ANSWER  7 ms, 0.9 MB
---- input ----
2
5 5
---- expected ----
10
---- got ----
11
---- stderr ----
debug: n=2
[cp] test 3: RUNTIME ERROR, exit 0xC00000FD: stack overflow (recursion too deep or a huge local array)  129 ms, 257.3 MB
[cp] 1/3 passed
\end{shell}

Compared with CPH's panel you get the reason for a crash, the peak memory of every test, and a warning when a
test takes longer than the problem's time limit. If the file has no CPH tests yet, F5 falls back to
\key{Ctrl+F5} behaviour. Output is compared token by token, so extra spaces or a missing final newline never count
as wrong.

\begin{center}
\begin{tabularx}{\linewidth}{>{\ttfamily}p{3.3cm}L}
\toprule
\thd{Exit code} & \thd{Meaning and usual cause} \\
\midrule
0xC00000FD & stack overflow: recursion deeper than the 256\,MB stack allows, or a huge local array (move it to a global or a vector) \\
0xC0000005 & access violation: out-of-bounds pointer or C-array access, dereferencing \texttt{end()} \\
0xC000001D & undefined behaviour trapped (only with \texttt{UB\_TRAP = True}, Section~\ref{sec:ubtrap}): signed overflow, bad index, bad shift \\
0xC0000094 & integer division or modulo by zero \\
0xC0000374 & heap corruption: writing past the end of an array \\
3 & \texttt{abort()}: a failed \texttt{assert}, or the checked STL caught a bug; its message is in stderr \\
\bottomrule
\end{tabularx}
\end{center}

\begin{fact}[The checked STL]
Every local build defines \texttt{\_GLIBCXX\_DEBUG}, so \cc{v[i]} with \cc{i} out of range, an invalid comparator in
\cc{sort}, or using an invalidated iterator stops the program with a message naming the problem instead of
silently producing garbage. That message appears in stderr, i.e.\ in CPH's Standard Error box or F5's output.
\end{fact}

\section{Ctrl+F5: typing input yourself}

\key{Ctrl+F5} compiles the file and runs it in a terminal where you type the input. Finish with \key{Ctrl+Z} then
\key{Enter} (end of file). When the program exits you get its exit code (decoded), CPU time and peak memory:

\begin{shell}
[cp] running sol.cpp  (type the input; Ctrl+Z then Enter sends EOF)
3
1 -2 5
4
[cp] exit 0 | cpu 0 ms | peak memory 0.9 MB | wall 4.21 s
\end{shell}

Binaries are cached in \file{.build/}: running again without editing the file skips compilation.

\section{Stress testing}\label{sec:stress}

When the samples pass but the judge disagrees, the bug hides in a case you have not thought of. A stress test
finds it automatically by comparing your solution with a slow but obviously correct brute force on thousands of
small random inputs.

\begin{figure}[H]
\centering
\begin{tikzpicture}[
  pr/.style={draw=cpblue, fill=cpblue!6, rounded corners=3pt, thick, minimum width=2.8cm, minimum height=1cm, align=center, font=\sffamily\small},
  dat/.style={draw=cpgray, fill=white, tape, minimum width=1.9cm, minimum height=0.9cm, align=center, font=\ttfamily\scriptsize},
  arr/.style={->, thick, cpgray}]
  \node[pr] (gen) at (0, 0) {\file{gen.cpp}\\\scriptsize seed $i$};
  \node[dat] (in) at (3.2, 0) {random\\input};
  \node[pr] (brute) at (6.6, 1.1) {\file{brute.cpp}\\\scriptsize slow, correct};
  \node[pr, draw=cporange, fill=cporange!6] (sol) at (6.6, -1.1) {your solution\\\scriptsize fast, suspect};
  \node[pr, draw=cpteal, fill=cpteal!6, diamond, aspect=1.6, minimum width=2.4cm] (cmp) at (10.2, 0) {same?};
  \draw[arr] (gen) -- (in);
  \draw[arr] (in) -- (brute);
  \draw[arr] (in) -- (sol);
  \draw[arr] (brute) -- (cmp);
  \draw[arr] (sol) -- (cmp);
  \draw[arr, cpgreen] (cmp.north) -- ++(0, 1.1) -| node[above, pos=0.25, font=\sffamily\scriptsize, text=cpgreen] {yes: next seed $i+1$} (gen.north);
  \draw[arr, cpred] (cmp.south) -- ++(0, -1.4) node[below, font=\sffamily\scriptsize, text=cpred, align=center] {no: print input, expected, got;\\copy the input to the clipboard};
\end{tikzpicture}
\caption{The stress loop run by \key{Ctrl+Alt+X}.}
\end{figure}

\subsection{How to use it}
\begin{enumerate}
  \item Edit \file{stress/gen.cpp} to print one random test in the problem's input format. The test number arrives as
        \cc{argv[1]} and seeds the generator, so a failure can be reproduced. Keep inputs \emph{tiny} ($n \le 8$):
        small failing cases are easy to understand.
  \item Edit \file{stress/brute.cpp}: the simplest correct solution you can write, however slow (try everything).
  \item Open your solution file and press \key{Ctrl+Alt+X}.
\end{enumerate}

The tester first runs \file{brute.cpp} on the problem's CPH samples; if the brute force fails a sample, it says so
and stops, because every later verdict would be meaningless. Then it runs 300 random tests (\texttt{-n} changes the
number when run by hand). On the first difference it prints the input, the brute force's output and yours, saves
the input to \file{stress/fail.in} and copies it to the clipboard. In CPH press \emph{New Testcase}, paste, type the
expected output, fix the bug and rerun.

\begin{shell}
[cp] brute.cpp passes all 1 CPH sample(s)
[cp] stress testing positives_only.cpp against brute.cpp (100 tests)
[cp] FAILED on test 1: wrong answer
---- input ----
2
-8 -1
---- expected (brute) ----
-9
---- got (yours) ----
0
[cp] input copied to the clipboard: in CPH press '+ New Testcase', paste it, and use the expected output above
\end{shell}

\subsection{Problems with several correct answers}
If any valid answer is accepted, comparing outputs is wrong. Add \file{stress/checker.cpp}: it receives the paths of
the input, the brute force's output and your output as \cc{argv[1..3]} and must exit with code 0 when your output
is valid. When the file exists, the tester uses it instead of comparing text.

\begin{tip}
Write the brute force and the generator \emph{before} you are desperate. Ten minutes spent on them usually beats an
hour of staring at code, and the same pair can be reused for the next attempt.
\end{tip}

\chapter{Compiler Flags and Speed}\label{ch:compiler}

Every local build (CPH, F5, Ctrl+F5, the stress tester, Code Runner) uses the same flags, defined once in
\file{tools/cp.py} and copied into CPH's settings by \emph{CP: Setup}:

\begin{shell}
-std=c++17 -O2 -DLOCAL -Wall -Wextra -Wshadow -D_GLIBCXX_DEBUG
-IC:/Users/kshit/Desktop/competitive/include -static -Wl,--stack,268435456
\end{shell}

\section{What each flag does}

\begin{center}
\begin{tabularx}{\linewidth}{>{\ttfamily\raggedright\arraybackslash}p{4.3cm}L}
\toprule
\thd{Flag} & \thd{Why it is there} \\
\midrule
-std=c++17 & The language standard you asked for. IntelliSense uses it too, so C++20 features are flagged while you type. \\
-O2 & The optimisation level judges use; timings are meaningful. \\
-DLOCAL & Defines \cc{LOCAL}: turns on \cc{dbg} and turns off USACO file I/O locally. \\
-Wall -Wextra -Wshadow & Warnings for common bugs: unused variables, sign comparisons, a local that hides another. \\
-D\_GLIBCXX\_DEBUG & The checked STL: out-of-range indexing, invalid comparators and bad iterators abort with a message. \\
-I\ldots/include & Where \file{debug.h} and the precompiled \file{bits/stdc++.h} live. \\
-static & The program carries its own C++ runtime. Other MinGW installs on your PATH (Git's, cpeditor's) can no longer
  make it load the wrong DLL and die silently. \\
-Wl,-{}-stack,268435456 & A 256\,MB stack, like Codeforces. Windows' default 2\,MB stack crashes a recursive DFS on a
  200\,000-vertex path; this was measured before and after. \\
\bottomrule
\end{tabularx}
\end{center}

CPH additionally passes \texttt{-D DEBUG -D CPH}, so \cc{#ifdef CPH} works if you ever want CPH-only code.

\section{The precompiled header}

Every compile used to spend most of its time parsing \file{bits/stdc++.h}, which pulls in the entire standard library.
\file{include/bits/} now holds a copy of that header plus \file{stdc++.h.gch}, the same header pre-parsed with exactly
the flags above. Because \texttt{-I} puts \file{include/} first in the search path, GCC finds the \texttt{.gch} and loads
it in a fraction of the time.

\begin{figure}[H]
\centering
\begin{tikzpicture}
\begin{axis}[xbar, width=12.5cm, height=4.6cm, bar width=9pt, xmin=0, xmax=1.7, enlarge y limits=0.45,
  xlabel={seconds per compile (warm cache)}, symbolic y coords={A\_Graph.cpp,empty template},
  ytick=data, y tick label style={font=\sffamily\small}, x tick label style={font=\sffamily\scriptsize},
  xlabel style={font=\sffamily\small}, legend style={font=\sffamily\scriptsize, at={(0.98,0.05)}, anchor=south east},
  nodes near coords, nodes near coords style={font=\sffamily\scriptsize}, every axis plot/.append style={draw=none}]
  \addplot[fill=cpgray!45] coordinates {(1.42,A\_Graph.cpp) (0.95,empty template)};
  \addplot[fill=cpteal] coordinates {(0.91,A\_Graph.cpp) (0.42,empty template)};
  \legend{old flags, new flags + precompiled header}
\end{axis}
\end{tikzpicture}
\caption{Measured CPH compile times. A cold cache (first compile after a reboot) went from about 3.5\,s to under 1\,s.}
\end{figure}

\begin{pitfall}[Keep CPH's flags in sync]
A precompiled header is only valid for the flags it was built with, and GCC does not always reject a mismatched one:
during setup a mismatch made GCC crash with an internal error instead of falling back. So never edit CPH's
\emph{Args} by hand. Change \texttt{COMPILE\_FLAGS} in \file{tools/cp.py} and run \emph{CP: Setup}; it rebuilds the
header and rewrites CPH's setting together. F5 and the stress tester also rebuild the header automatically whenever
the compiler or the flags change.
\end{pitfall}

\section{The optional UB trap}\label{sec:ubtrap}

\file{tools/cp.py} has a switch:

\begin{code}
UB_TRAP = False
\end{code}

Set to \cc{True}, it adds \texttt{-fsanitize=undefined -fsanitize-undefined-trap-on-error}: signed overflow,
out-of-bounds C-array indexing and invalid shifts crash with exit code \texttt{0xC000001D} instead of silently
producing a wrong answer. It was tested and works with static linking, but costs about 0.35\,s per compile, which
cancels most of the precompiled header's gain, so it is off by default. With \cc{#define int long long} and the
checked STL most of what it catches is already covered. Turn it on (then run \emph{CP: Setup}) when you suspect
overflow.

\section{Time limits and the checked STL}

The checked STL verifies every vector access, which can make STL-heavy code several times slower locally than on
the judge. If CPH or F5 reports a time close to the limit, that is usually not a real TLE. To measure true speed,
temporarily remove \texttt{-D\_GLIBCXX\_DEBUG} from \texttt{COMPILE\_FLAGS} and run \emph{CP: Setup}; put it back afterwards.

\chapter{The Template Library}\label{ch:library}

\file{library/} holds 174 templates plus the full file template, 175 snippets in all, covering the USACO Bronze
to Platinum syllabus and the advanced topics beyond it. The second edition of the handbook added 45 of them (discrete
logarithms, polynomial tools, FWHT, Du's sieve and prime counting, Hungarian and blossom matching, Gomory--Hu and
Stoer--Wagner cuts, the eertree, Lyndon factorization, CDQ, Mo on trees and with updates, persistent treaps, Euler tour
trees, half-planes, rotating calipers and more). Every one of them:

\begin{itemize}
  \item compiles on its own as strict C++17 with zero warnings (\texttt{-Wall -Wextra -Wshadow -Wpedantic}), both as
        written and under \cc{#define int long long};
  \item was checked against a brute-force implementation on hundreds to thousands of random inputs (and planted bugs
        were confirmed to be caught);
  \item avoids every name the template defines (\cc{INF}, \cc{MOD}, \cc{sz}, \cc{all}, \cc{rep}, \cc{dbg},
        \ldots) and uses no \cc{__int128}, so it also works on 32-bit judges.
\end{itemize}

\section{Three ways to insert a template}

\begin{enumerate}
  \item \textbf{Chords.} \key{Ctrl+Alt+T} (core), \key{Ctrl+Alt+Y} (advanced) or \key{Ctrl+Alt+U} (specialist
        techniques), release, then one key. The code is pasted at the cursor and re-indented to match it. The maps below
        show every chord.
  \item \textbf{Prefixes.} Every template has a prefix starting with \texttt{cp}: type it and press \key{Tab}
        (\texttt{cpdsu}, \texttt{cpseg}, \texttt{cphld}). Typing just \texttt{cp} lists all of them with their descriptions.
  \item \textbf{Picker.} \key{Ctrl+Alt+T} then \key{Space} opens VS Code's snippet list; type any part of a name.
\end{enumerate}

\section{Where to paste}

Each template is one of two kinds, shown on its card in Chapter~\ref{ch:reference} and in \file{CHEATSHEET.md}:

\begin{itemize}
  \item \textbf{above solve()} (\emph{global}): a \cc{struct} or functions, such as \cc{DSU}, \cc{SegTree},
        \cc{dijkstra}. Put the cursor on an empty line above \cc{void solve()} and paste; then use it inside
        \cc{solve()}.
  \item \textbf{inside solve()}: a block of statements that uses variables you already have, such as the DFS or grid
        BFS. The first line of the card says which variables it needs (\texttt{needs n, adj}). Paste it where those
        variables exist, then fill in the marked lines.
\end{itemize}

A few templates build on another one and say so (\texttt{Needs cppoint}, \texttt{Needs cpdinic}, \texttt{Needs cpdsurb}):
paste the dependency first.

\section{Conventions used everywhere}

\begin{center}
\begin{tabularx}{\linewidth}{p{4cm}L}
\toprule
\thd{Convention} & \thd{Meaning} \\
\midrule
0-indexed & vertices $0..n-1$, array positions $0..n-1$. \texttt{cpgraph} converts 1-indexed input for you. \\
half-open ranges & every range query is on $[l, r)$: \cc{st.query(0, n)} is the whole array, \cc{query(l, l)} is empty. \\
unreachable & shortest-path templates return \cc{LLONG_MAX} (or the named constant they define) for unreachable vertices. \\
empty result = impossible & e.g.\ \cc{topo_sort} returns an empty vector on a cycle, \cc{TwoSat::solve} when unsatisfiable. \\
\texttt{MOD} & modular templates use the template's \cc{MOD}. \\
edit points & lines to adapt are marked with a comment (\texttt{// transition}, \texttt{// replace with the real check}). \\
\bottomrule
\end{tabularx}
\end{center}

\section{Keyboard maps}

Orange keys have a template; the short name under the key is its prefix without \texttt{cp}. The complete list with
descriptions is Chapter~\ref{ch:reference}.

% generated by tools/cp.py setup -- do not edit
\begin{figure}[H]\centering
\begin{tikzpicture}[kb/.style={draw=cpgray!45, fill=white, rounded corners=2pt, minimum width=11.6mm, minimum height=10.4mm, align=center, font=\sffamily\fontsize{5.6}{6.4}\selectfont, inner sep=0.6pt, text=cpgray!70}, kbused/.style={fill=cporange!10, draw=cporange!70, text=black}]
\node[kb, kbused] at (0.00,0.00) {\textbf{\color{cpblue}\textasciigrave{}}\\xortrie};
\node[kb, kbused] at (1.24,0.00) {\textbf{\color{cpblue}1}\\prefix};
\node[kb, kbused] at (2.48,0.00) {\textbf{\color{cpblue}2}\\bipartite};
\node[kb, kbused] at (3.72,0.00) {\textbf{\color{cpblue}3}\\diff};
\node[kb, kbused] at (4.96,0.00) {\textbf{\color{cpblue}4}\\diff2d};
\node[kb, kbused] at (6.20,0.00) {\textbf{\color{cpblue}5}\\twoptr};
\node[kb, kbused] at (7.44,0.00) {\textbf{\color{cpblue}6}\\window};
\node[kb, kbused] at (8.68,0.00) {\textbf{\color{cpblue}7}\\monostack};
\node[kb, kbused] at (9.92,0.00) {\textbf{\color{cpblue}8}\\subsets};
\node[kb, kbused] at (11.16,0.00) {\textbf{\color{cpblue}9}\\perm};
\node[kb, kbused] at (12.40,0.00) {\textbf{\color{cpblue}0}\\01bfs};
\node[kb, kbused] at (13.64,0.00) {\textbf{\color{cpblue}-}\\rangedp};
\node[kb, kbused] at (14.88,0.00) {\textbf{\color{cpblue}=}\\digit};
\node[kb, kbused] at (0.62,-1.12) {\textbf{\color{cpblue}Q}\\ternary};
\node[kb, kbused] at (1.86,-1.12) {\textbf{\color{cpblue}W}\\floyd};
\node[kb, kbused] at (3.10,-1.12) {\textbf{\color{cpblue}E}\\tree};
\node[kb, kbused] at (4.34,-1.12) {\textbf{\color{cpblue}R}\\sparse};
\node[kb, kbused] at (5.58,-1.12) {\textbf{\color{cpblue}T}\\template};
\node[kb, kbused] at (6.82,-1.12) {\textbf{\color{cpblue}Y}\\bits};
\node[kb, kbused] at (8.06,-1.12) {\textbf{\color{cpblue}U}\\dsu};
\node[kb, kbused] at (9.30,-1.12) {\textbf{\color{cpblue}I}\\bit};
\node[kb, kbused] at (10.54,-1.12) {\textbf{\color{cpblue}O}\\topo};
\node[kb, kbused] at (11.78,-1.12) {\textbf{\color{cpblue}P}\\sieve};
\node[kb, kbused] at (13.02,-1.12) {\textbf{\color{cpblue}[}\\lcs};
\node[kb, kbused] at (14.26,-1.12) {\textbf{\color{cpblue}]}\\tsp};
\node[kb, kbused] at (0.93,-2.24) {\textbf{\color{cpblue}A}\\bsearch};
\node[kb, kbused] at (2.17,-2.24) {\textbf{\color{cpblue}S}\\seg};
\node[kb, kbused] at (3.41,-2.24) {\textbf{\color{cpblue}D}\\dfs};
\node[kb, kbused] at (4.65,-2.24) {\textbf{\color{cpblue}F}\\flood};
\node[kb, kbused] at (5.89,-2.24) {\textbf{\color{cpblue}G}\\gridbfs};
\node[kb, kbused] at (7.13,-2.24) {\textbf{\color{cpblue}H}\\strhash};
\node[kb, kbused] at (8.37,-2.24) {\textbf{\color{cpblue}J}\\dijkstra};
\node[kb, kbused] at (9.61,-2.24) {\textbf{\color{cpblue}K}\\kruskal};
\node[kb, kbused] at (10.85,-2.24) {\textbf{\color{cpblue}L}\\lca};
\node[kb, kbused] at (12.09,-2.24) {\textbf{\color{cpblue};}\\treedp};
\node[kb, kbused] at (13.33,-2.24) {\textbf{\color{cpblue}\textquotesingle{}}\\reroot};
\node[kb, kbused] at (1.55,-3.36) {\textbf{\color{cpblue}Z}\\z};
\node[kb, kbused] at (2.79,-3.36) {\textbf{\color{cpblue}X}\\extgcd};
\node[kb, kbused] at (4.03,-3.36) {\textbf{\color{cpblue}C}\\scc};
\node[kb, kbused] at (5.27,-3.36) {\textbf{\color{cpblue}V}\\read};
\node[kb, kbused] at (6.51,-3.36) {\textbf{\color{cpblue}B}\\bfs};
\node[kb, kbused] at (7.75,-3.36) {\textbf{\color{cpblue}N}\\ncr};
\node[kb, kbused] at (8.99,-3.36) {\textbf{\color{cpblue}M}\\modpow};
\node[kb, kbused] at (10.23,-3.36) {\textbf{\color{cpblue},}\\funcgraph};
\node[kb, kbused] at (11.47,-3.36) {\textbf{\color{cpblue}.}\\jump};
\node[kb, kbused] at (12.71,-3.36) {\textbf{\color{cpblue}/}\\cycle};
\end{tikzpicture}
\caption{Ctrl+Alt+T, then a key no Shift (46 templates)}
\end{figure}
\begin{figure}[H]\centering
\begin{tikzpicture}[kb/.style={draw=cpgray!45, fill=white, rounded corners=2pt, minimum width=11.6mm, minimum height=10.4mm, align=center, font=\sffamily\fontsize{5.6}{6.4}\selectfont, inner sep=0.6pt, text=cpgray!70}, kbused/.style={fill=cporange!10, draw=cporange!70, text=black}]
\node[kb] at (0.00,0.00) {\textbf{\color{cpblue}\textasciigrave{}}};
\node[kb, kbused] at (1.24,0.00) {\textbf{\color{cpblue}1}\\bitsetsum};
\node[kb, kbused] at (2.48,0.00) {\textbf{\color{cpblue}2}\\divsieve};
\node[kb, kbused] at (3.72,0.00) {\textbf{\color{cpblue}3}\\point};
\node[kb, kbused] at (4.96,0.00) {\textbf{\color{cpblue}4}\\rect};
\node[kb, kbused] at (6.20,0.00) {\textbf{\color{cpblue}5}\\prim};
\node[kb] at (7.44,0.00) {\textbf{\color{cpblue}6}};
\node[kb] at (8.68,0.00) {\textbf{\color{cpblue}7}};
\node[kb] at (9.92,0.00) {\textbf{\color{cpblue}8}};
\node[kb] at (11.16,0.00) {\textbf{\color{cpblue}9}};
\node[kb] at (12.40,0.00) {\textbf{\color{cpblue}0}};
\node[kb] at (13.64,0.00) {\textbf{\color{cpblue}-}};
\node[kb] at (14.88,0.00) {\textbf{\color{cpblue}=}};
\node[kb, kbused] at (0.62,-1.12) {\textbf{\color{cpblue}Q}\\knap};
\node[kb, kbused] at (1.86,-1.12) {\textbf{\color{cpblue}W}\\bellman};
\node[kb, kbused] at (3.10,-1.12) {\textbf{\color{cpblue}E}\\graphw};
\node[kb, kbused] at (4.34,-1.12) {\textbf{\color{cpblue}R}\\rng};
\node[kb, kbused] at (5.58,-1.12) {\textbf{\color{cpblue}T}\\trie};
\node[kb, kbused] at (6.82,-1.12) {\textbf{\color{cpblue}Y}\\backtrack};
\node[kb, kbused] at (8.06,-1.12) {\textbf{\color{cpblue}U}\\mitm};
\node[kb, kbused] at (9.30,-1.12) {\textbf{\color{cpblue}I}\\interacti.};
\node[kb, kbused] at (10.54,-1.12) {\textbf{\color{cpblue}O}\\oset};
\node[kb, kbused] at (11.78,-1.12) {\textbf{\color{cpblue}P}\\factor};
\node[kb] at (13.02,-1.12) {\textbf{\color{cpblue}[}};
\node[kb] at (14.26,-1.12) {\textbf{\color{cpblue}]}};
\node[kb, kbused] at (0.93,-2.24) {\textbf{\color{cpblue}A}\\prefix2d};
\node[kb, kbused] at (2.17,-2.24) {\textbf{\color{cpblue}S}\\lazy};
\node[kb, kbused] at (3.41,-2.24) {\textbf{\color{cpblue}D}\\diameter};
\node[kb, kbused] at (4.65,-2.24) {\textbf{\color{cpblue}F}\\bsreal};
\node[kb, kbused] at (5.89,-2.24) {\textbf{\color{cpblue}G}\\graph};
\node[kb, kbused] at (7.13,-2.24) {\textbf{\color{cpblue}H}\\hash};
\node[kb, kbused] at (8.37,-2.24) {\textbf{\color{cpblue}J}\\argsort};
\node[kb, kbused] at (9.61,-2.24) {\textbf{\color{cpblue}K}\\kmp};
\node[kb, kbused] at (10.85,-2.24) {\textbf{\color{cpblue}L}\\lis};
\node[kb] at (12.09,-2.24) {\textbf{\color{cpblue};}};
\node[kb] at (13.33,-2.24) {\textbf{\color{cpblue}\textquotesingle{}}};
\node[kb, kbused] at (1.55,-3.36) {\textbf{\color{cpblue}Z}\\gridpaths};
\node[kb, kbused] at (2.79,-3.36) {\textbf{\color{cpblue}X}\\matrix};
\node[kb, kbused] at (4.03,-3.36) {\textbf{\color{cpblue}C}\\compress};
\node[kb, kbused] at (5.27,-3.36) {\textbf{\color{cpblue}V}\\hull};
\node[kb, kbused] at (6.51,-3.36) {\textbf{\color{cpblue}B}\\bridges};
\node[kb, kbused] at (7.75,-3.36) {\textbf{\color{cpblue}N}\\intervals};
\node[kb, kbused] at (8.99,-3.36) {\textbf{\color{cpblue}M}\\mint};
\node[kb] at (10.23,-3.36) {\textbf{\color{cpblue},}};
\node[kb] at (11.47,-3.36) {\textbf{\color{cpblue}.}};
\node[kb] at (12.71,-3.36) {\textbf{\color{cpblue}/}};
\end{tikzpicture}
\caption{Ctrl+Alt+T, then a key with Shift (31 templates)}
\end{figure}
\begin{figure}[H]\centering
\begin{tikzpicture}[kb/.style={draw=cpgray!45, fill=white, rounded corners=2pt, minimum width=11.6mm, minimum height=10.4mm, align=center, font=\sffamily\fontsize{5.6}{6.4}\selectfont, inner sep=0.6pt, text=cpgray!70}, kbused/.style={fill=cporange!10, draw=cporange!70, text=black}]
\node[kb] at (0.00,0.00) {\textbf{\color{cpblue}\textasciigrave{}}};
\node[kb, kbused] at (1.24,0.00) {\textbf{\color{cpblue}1}\\crt};
\node[kb, kbused] at (2.48,0.00) {\textbf{\color{cpblue}2}\\2sat};
\node[kb, kbused] at (3.72,0.00) {\textbf{\color{cpblue}3}\\incexc};
\node[kb, kbused] at (4.96,0.00) {\textbf{\color{cpblue}4}\\sa};
\node[kb, kbused] at (6.20,0.00) {\textbf{\color{cpblue}5}\\sam};
\node[kb, kbused] at (7.44,0.00) {\textbf{\color{cpblue}6}\\manacher};
\node[kb, kbused] at (8.68,0.00) {\textbf{\color{cpblue}7}\\seginter};
\node[kb, kbused] at (9.92,0.00) {\textbf{\color{cpblue}8}\\polygon};
\node[kb, kbused] at (11.16,0.00) {\textbf{\color{cpblue}9}\\closest};
\node[kb] at (12.40,0.00) {\textbf{\color{cpblue}0}};
\node[kb] at (13.64,0.00) {\textbf{\color{cpblue}-}};
\node[kb] at (14.88,0.00) {\textbf{\color{cpblue}=}};
\node[kb, kbused] at (0.62,-1.12) {\textbf{\color{cpblue}Q}\\dnc};
\node[kb, kbused] at (1.86,-1.12) {\textbf{\color{cpblue}W}\\wavelet};
\node[kb, kbused] at (3.10,-1.12) {\textbf{\color{cpblue}E}\\euler};
\node[kb, kbused] at (4.34,-1.12) {\textbf{\color{cpblue}R}\\dsurb};
\node[kb, kbused] at (5.58,-1.12) {\textbf{\color{cpblue}T}\\treap};
\node[kb, kbused] at (6.82,-1.12) {\textbf{\color{cpblue}Y}\\dinic};
\node[kb, kbused] at (8.06,-1.12) {\textbf{\color{cpblue}U}\\sos};
\node[kb, kbused] at (9.30,-1.12) {\textbf{\color{cpblue}I}\\bitrange};
\node[kb, kbused] at (10.54,-1.12) {\textbf{\color{cpblue}O}\\mo};
\node[kb, kbused] at (11.78,-1.12) {\textbf{\color{cpblue}P}\\persist};
\node[kb] at (13.02,-1.12) {\textbf{\color{cpblue}[}};
\node[kb] at (14.26,-1.12) {\textbf{\color{cpblue}]}};
\node[kb, kbused] at (0.93,-2.24) {\textbf{\color{cpblue}A}\\lazyassign};
\node[kb, kbused] at (2.17,-2.24) {\textbf{\color{cpblue}S}\\beats};
\node[kb, kbused] at (3.41,-2.24) {\textbf{\color{cpblue}D}\\dynseg};
\node[kb, kbused] at (4.65,-2.24) {\textbf{\color{cpblue}F}\\fft};
\node[kb, kbused] at (5.89,-2.24) {\textbf{\color{cpblue}G}\\grundy};
\node[kb, kbused] at (7.13,-2.24) {\textbf{\color{cpblue}H}\\hld};
\node[kb] at (8.37,-2.24) {\textbf{\color{cpblue}J}};
\node[kb, kbused] at (9.61,-2.24) {\textbf{\color{cpblue}K}\\krt};
\node[kb, kbused] at (10.85,-2.24) {\textbf{\color{cpblue}L}\\lineconta.};
\node[kb] at (12.09,-2.24) {\textbf{\color{cpblue};}};
\node[kb] at (13.33,-2.24) {\textbf{\color{cpblue}\textquotesingle{}}};
\node[kb, kbused] at (1.55,-3.36) {\textbf{\color{cpblue}Z}\\monocht};
\node[kb, kbused] at (2.79,-3.36) {\textbf{\color{cpblue}X}\\pollard};
\node[kb, kbused] at (4.03,-3.36) {\textbf{\color{cpblue}C}\\centroid};
\node[kb, kbused] at (5.27,-3.36) {\textbf{\color{cpblue}V}\\virtual};
\node[kb, kbused] at (6.51,-3.36) {\textbf{\color{cpblue}B}\\2ecc};
\node[kb, kbused] at (7.75,-3.36) {\textbf{\color{cpblue}N}\\segnode};
\node[kb, kbused] at (8.99,-3.36) {\textbf{\color{cpblue}M}\\small2lar.};
\node[kb] at (10.23,-3.36) {\textbf{\color{cpblue},}};
\node[kb] at (11.47,-3.36) {\textbf{\color{cpblue}.}};
\node[kb] at (12.71,-3.36) {\textbf{\color{cpblue}/}};
\end{tikzpicture}
\caption{Ctrl+Alt+Y, then a key no Shift (34 templates)}
\end{figure}
\begin{figure}[H]\centering
\begin{tikzpicture}[kb/.style={draw=cpgray!45, fill=white, rounded corners=2pt, minimum width=11.6mm, minimum height=10.4mm, align=center, font=\sffamily\fontsize{5.6}{6.4}\selectfont, inner sep=0.6pt, text=cpgray!70}, kbused/.style={fill=cporange!10, draw=cporange!70, text=black}]
\node[kb] at (0.00,0.00) {\textbf{\color{cpblue}\textasciigrave{}}};
\node[kb] at (1.24,0.00) {\textbf{\color{cpblue}1}};
\node[kb] at (2.48,0.00) {\textbf{\color{cpblue}2}};
\node[kb] at (3.72,0.00) {\textbf{\color{cpblue}3}};
\node[kb] at (4.96,0.00) {\textbf{\color{cpblue}4}};
\node[kb] at (6.20,0.00) {\textbf{\color{cpblue}5}};
\node[kb] at (7.44,0.00) {\textbf{\color{cpblue}6}};
\node[kb] at (8.68,0.00) {\textbf{\color{cpblue}7}};
\node[kb] at (9.92,0.00) {\textbf{\color{cpblue}8}};
\node[kb] at (11.16,0.00) {\textbf{\color{cpblue}9}};
\node[kb] at (12.40,0.00) {\textbf{\color{cpblue}0}};
\node[kb] at (13.64,0.00) {\textbf{\color{cpblue}-}};
\node[kb] at (14.88,0.00) {\textbf{\color{cpblue}=}};
\node[kb, kbused] at (0.62,-1.12) {\textbf{\color{cpblue}Q}\\knuth};
\node[kb] at (1.86,-1.12) {\textbf{\color{cpblue}W}};
\node[kb, kbused] at (3.10,-1.12) {\textbf{\color{cpblue}E}\\phi};
\node[kb, kbused] at (4.34,-1.12) {\textbf{\color{cpblue}R}\\offlinecc};
\node[kb, kbused] at (5.58,-1.12) {\textbf{\color{cpblue}T}\\ntt};
\node[kb, kbused] at (6.82,-1.12) {\textbf{\color{cpblue}Y}\\mcmf};
\node[kb] at (8.06,-1.12) {\textbf{\color{cpblue}U}};
\node[kb, kbused] at (9.30,-1.12) {\textbf{\color{cpblue}I}\\bit2d};
\node[kb] at (10.54,-1.12) {\textbf{\color{cpblue}O}};
\node[kb, kbused] at (11.78,-1.12) {\textbf{\color{cpblue}P}\\profile};
\node[kb] at (13.02,-1.12) {\textbf{\color{cpblue}[}};
\node[kb] at (14.26,-1.12) {\textbf{\color{cpblue}]}};
\node[kb, kbused] at (0.93,-2.24) {\textbf{\color{cpblue}A}\\aho};
\node[kb, kbused] at (2.17,-2.24) {\textbf{\color{cpblue}S}\\slope};
\node[kb, kbused] at (3.41,-2.24) {\textbf{\color{cpblue}D}\\distinct};
\node[kb, kbused] at (4.65,-2.24) {\textbf{\color{cpblue}F}\\floorbloc.};
\node[kb, kbused] at (5.89,-2.24) {\textbf{\color{cpblue}G}\\flowlb};
\node[kb, kbused] at (7.13,-2.24) {\textbf{\color{cpblue}H}\\hopcroft};
\node[kb] at (8.37,-2.24) {\textbf{\color{cpblue}J}};
\node[kb, kbused] at (9.61,-2.24) {\textbf{\color{cpblue}K}\\treeknap};
\node[kb, kbused] at (10.85,-2.24) {\textbf{\color{cpblue}L}\\lichao};
\node[kb] at (12.09,-2.24) {\textbf{\color{cpblue};}};
\node[kb] at (13.33,-2.24) {\textbf{\color{cpblue}\textquotesingle{}}};
\node[kb] at (1.55,-3.36) {\textbf{\color{cpblue}Z}};
\node[kb, kbused] at (2.79,-3.36) {\textbf{\color{cpblue}X}\\xorbasis};
\node[kb, kbused] at (4.03,-3.36) {\textbf{\color{cpblue}C}\\mincount};
\node[kb] at (5.27,-3.36) {\textbf{\color{cpblue}V}};
\node[kb, kbused] at (6.51,-3.36) {\textbf{\color{cpblue}B}\\bcc};
\node[kb, kbused] at (7.75,-3.36) {\textbf{\color{cpblue}N}\\pascal};
\node[kb, kbused] at (8.99,-3.36) {\textbf{\color{cpblue}M}\\mobius};
\node[kb] at (10.23,-3.36) {\textbf{\color{cpblue},}};
\node[kb] at (11.47,-3.36) {\textbf{\color{cpblue}.}};
\node[kb] at (12.71,-3.36) {\textbf{\color{cpblue}/}};
\end{tikzpicture}
\caption{Ctrl+Alt+Y, then a key with Shift (20 templates)}
\end{figure}
\begin{figure}[H]\centering
\begin{tikzpicture}[kb/.style={draw=cpgray!45, fill=white, rounded corners=2pt, minimum width=11.6mm, minimum height=10.4mm, align=center, font=\sffamily\fontsize{5.6}{6.4}\selectfont, inner sep=0.6pt, text=cpgray!70}, kbused/.style={fill=cporange!10, draw=cporange!70, text=black}]
\node[kb] at (0.00,0.00) {\textbf{\color{cpblue}\textasciigrave{}}};
\node[kb] at (1.24,0.00) {\textbf{\color{cpblue}1}};
\node[kb, kbused] at (2.48,0.00) {\textbf{\color{cpblue}2}\\bibfs};
\node[kb, kbused] at (3.72,0.00) {\textbf{\color{cpblue}3}\\cdq};
\node[kb, kbused] at (4.96,0.00) {\textbf{\color{cpblue}4}\\halfplane};
\node[kb, kbused] at (6.20,0.00) {\textbf{\color{cpblue}5}\\calipers};
\node[kb] at (7.44,0.00) {\textbf{\color{cpblue}6}};
\node[kb] at (8.68,0.00) {\textbf{\color{cpblue}7}};
\node[kb] at (9.92,0.00) {\textbf{\color{cpblue}8}};
\node[kb, kbused] at (11.16,0.00) {\textbf{\color{cpblue}9}\\debruijn};
\node[kb, kbused] at (12.40,0.00) {\textbf{\color{cpblue}0}\\gray};
\node[kb] at (13.64,0.00) {\textbf{\color{cpblue}-}};
\node[kb] at (14.88,0.00) {\textbf{\color{cpblue}=}};
\node[kb, kbused] at (0.62,-1.12) {\textbf{\color{cpblue}Q}\\sqrtmod};
\node[kb, kbused] at (1.86,-1.12) {\textbf{\color{cpblue}W}\\fwht};
\node[kb, kbused] at (3.10,-1.12) {\textbf{\color{cpblue}E}\\ett};
\node[kb, kbused] at (4.34,-1.12) {\textbf{\color{cpblue}R}\\ptreap};
\node[kb, kbused] at (5.58,-1.12) {\textbf{\color{cpblue}T}\\disjoints.};
\node[kb, kbused] at (6.82,-1.12) {\textbf{\color{cpblue}Y}\\lyndon};
\node[kb, kbused] at (8.06,-1.12) {\textbf{\color{cpblue}U}\\prufer};
\node[kb, kbused] at (9.30,-1.12) {\textbf{\color{cpblue}I}\\lagrange};
\node[kb, kbused] at (10.54,-1.12) {\textbf{\color{cpblue}O}\\motree};
\node[kb, kbused] at (11.78,-1.12) {\textbf{\color{cpblue}P}\\primecount};
\node[kb] at (13.02,-1.12) {\textbf{\color{cpblue}[}};
\node[kb] at (14.26,-1.12) {\textbf{\color{cpblue}]}};
\node[kb, kbused] at (0.93,-2.24) {\textbf{\color{cpblue}A}\\gomoryhu};
\node[kb, kbused] at (2.17,-2.24) {\textbf{\color{cpblue}S}\\stirling};
\node[kb, kbused] at (3.41,-2.24) {\textbf{\color{cpblue}D}\\dusieve};
\node[kb, kbused] at (4.65,-2.24) {\textbf{\color{cpblue}F}\\cyclefind};
\node[kb, kbused] at (5.89,-2.24) {\textbf{\color{cpblue}G}\\primroot};
\node[kb, kbused] at (7.13,-2.24) {\textbf{\color{cpblue}H}\\treehash};
\node[kb, kbused] at (8.37,-2.24) {\textbf{\color{cpblue}J}\\jumpptr};
\node[kb, kbused] at (9.61,-2.24) {\textbf{\color{cpblue}K}\\kitamasa};
\node[kb, kbused] at (10.85,-2.24) {\textbf{\color{cpblue}L}\\dlog};
\node[kb] at (12.09,-2.24) {\textbf{\color{cpblue};}};
\node[kb] at (13.33,-2.24) {\textbf{\color{cpblue}\textquotesingle{}}};
\node[kb, kbused] at (1.55,-3.36) {\textbf{\color{cpblue}Z}\\anneal};
\node[kb, kbused] at (2.79,-3.36) {\textbf{\color{cpblue}X}\\xorhash};
\node[kb, kbused] at (4.03,-3.36) {\textbf{\color{cpblue}C}\\cartesian};
\node[kb, kbused] at (5.27,-3.36) {\textbf{\color{cpblue}V}\\konig};
\node[kb, kbused] at (6.51,-3.36) {\textbf{\color{cpblue}B}\\sternbroc.};
\node[kb, kbused] at (7.75,-3.36) {\textbf{\color{cpblue}N}\\gauss};
\node[kb, kbused] at (8.99,-3.36) {\textbf{\color{cpblue}M}\\monominima};
\node[kb] at (10.23,-3.36) {\textbf{\color{cpblue},}};
\node[kb] at (11.47,-3.36) {\textbf{\color{cpblue}.}};
\node[kb] at (12.71,-3.36) {\textbf{\color{cpblue}/}};
\end{tikzpicture}
\caption{Ctrl+Alt+U, then a key no Shift (32 templates)}
\end{figure}
\begin{figure}[H]\centering
\begin{tikzpicture}[kb/.style={draw=cpgray!45, fill=white, rounded corners=2pt, minimum width=11.6mm, minimum height=10.4mm, align=center, font=\sffamily\fontsize{5.6}{6.4}\selectfont, inner sep=0.6pt, text=cpgray!70}, kbused/.style={fill=cporange!10, draw=cporange!70, text=black}]
\node[kb] at (0.00,0.00) {\textbf{\color{cpblue}\textasciigrave{}}};
\node[kb] at (1.24,0.00) {\textbf{\color{cpblue}1}};
\node[kb] at (2.48,0.00) {\textbf{\color{cpblue}2}};
\node[kb] at (3.72,0.00) {\textbf{\color{cpblue}3}};
\node[kb] at (4.96,0.00) {\textbf{\color{cpblue}4}};
\node[kb] at (6.20,0.00) {\textbf{\color{cpblue}5}};
\node[kb] at (7.44,0.00) {\textbf{\color{cpblue}6}};
\node[kb] at (8.68,0.00) {\textbf{\color{cpblue}7}};
\node[kb] at (9.92,0.00) {\textbf{\color{cpblue}8}};
\node[kb] at (11.16,0.00) {\textbf{\color{cpblue}9}};
\node[kb] at (12.40,0.00) {\textbf{\color{cpblue}0}};
\node[kb] at (13.64,0.00) {\textbf{\color{cpblue}-}};
\node[kb] at (14.88,0.00) {\textbf{\color{cpblue}=}};
\node[kb] at (0.62,-1.12) {\textbf{\color{cpblue}Q}};
\node[kb] at (1.86,-1.12) {\textbf{\color{cpblue}W}};
\node[kb, kbused] at (3.10,-1.12) {\textbf{\color{cpblue}E}\\eertree};
\node[kb, kbused] at (4.34,-1.12) {\textbf{\color{cpblue}R}\\radix};
\node[kb, kbused] at (5.58,-1.12) {\textbf{\color{cpblue}T}\\sack};
\node[kb] at (6.82,-1.12) {\textbf{\color{cpblue}Y}};
\node[kb] at (8.06,-1.12) {\textbf{\color{cpblue}U}};
\node[kb, kbused] at (9.30,-1.12) {\textbf{\color{cpblue}I}\\diffcons};
\node[kb, kbused] at (10.54,-1.12) {\textbf{\color{cpblue}O}\\moupdates};
\node[kb, kbused] at (11.78,-1.12) {\textbf{\color{cpblue}P}\\poly};
\node[kb] at (13.02,-1.12) {\textbf{\color{cpblue}[}};
\node[kb] at (14.26,-1.12) {\textbf{\color{cpblue}]}};
\node[kb, kbused] at (0.93,-2.24) {\textbf{\color{cpblue}A}\\stoerwagn.};
\node[kb, kbused] at (2.17,-2.24) {\textbf{\color{cpblue}S}\\subsetconv};
\node[kb, kbused] at (3.41,-2.24) {\textbf{\color{cpblue}D}\\det};
\node[kb] at (4.65,-2.24) {\textbf{\color{cpblue}F}};
\node[kb, kbused] at (5.89,-2.24) {\textbf{\color{cpblue}G}\\blossom};
\node[kb, kbused] at (7.13,-2.24) {\textbf{\color{cpblue}H}\\hungarian};
\node[kb] at (8.37,-2.24) {\textbf{\color{cpblue}J}};
\node[kb] at (9.61,-2.24) {\textbf{\color{cpblue}K}};
\node[kb, kbused] at (10.85,-2.24) {\textbf{\color{cpblue}L}\\offlinelca};
\node[kb] at (12.09,-2.24) {\textbf{\color{cpblue};}};
\node[kb] at (13.33,-2.24) {\textbf{\color{cpblue}\textquotesingle{}}};
\node[kb] at (1.55,-3.36) {\textbf{\color{cpblue}Z}};
\node[kb] at (2.79,-3.36) {\textbf{\color{cpblue}X}};
\node[kb, kbused] at (4.03,-3.36) {\textbf{\color{cpblue}C}\\minplus};
\node[kb] at (5.27,-3.36) {\textbf{\color{cpblue}V}};
\node[kb, kbused] at (6.51,-3.36) {\textbf{\color{cpblue}B}\\bm};
\node[kb, kbused] at (7.75,-3.36) {\textbf{\color{cpblue}N}\\longpath};
\node[kb, kbused] at (8.99,-3.36) {\textbf{\color{cpblue}M}\\mincircle};
\node[kb] at (10.23,-3.36) {\textbf{\color{cpblue},}};
\node[kb] at (11.47,-3.36) {\textbf{\color{cpblue}.}};
\node[kb] at (12.71,-3.36) {\textbf{\color{cpblue}/}};
\end{tikzpicture}
\caption{Ctrl+Alt+U, then a key with Shift (16 templates)}
\end{figure}


\section{Categories}

\begin{center}
\begin{tabularx}{\linewidth}{p{3.4cm}L}
\toprule
\thd{Folder} & \thd{Contents} \\
\midrule
\file{basics/} & input readers, prefix sums, difference arrays, two pointers, sliding window, monotonic stack, binary and
  ternary search, subsets, permutations, backtracking, argsort, intervals, bit helpers, meet in the middle,
  coordinate compression \\
\file{graph/} & DFS, BFS, flood fill, grid BFS, Dijkstra, 0-1 BFS, Bellman--Ford, Floyd--Warshall, topological sort,
  DSU (+ rollback), Kruskal, Prim, SCC, 2-SAT, bipartite check, bridges and cut vertices, 2-edge-connected components,
  biconnected components, Euler paths, cycle finding, functional graphs, k-th successor, offline connectivity \\
\file{tree/} & tree DFS / Euler tour, diameter, LCA, tree DP, rerooting, small-to-large, HLD, centroid decomposition,
  virtual tree, Kruskal reconstruction tree, tree knapsack \\
\file{ds/} & segment trees (point, lazy add, lazy assign, custom node with binary search, min + count, dynamic,
  persistent, beats), Fenwick trees (point, range, 2D), sparse table, ordered set, anti-hack hash, wavelet tree,
  treap, Mo's algorithm, offline distinct count \\
\file{dp/} & knapsacks, bitset subset sum, grid paths, LIS, LCS / edit distance, bitmask (TSP), range DP, digit DP,
  SOS DP, divide-and-conquer and Knuth optimisations, broken profile, LineContainer, Li Chao tree, slope trick \\
\file{math/} & modular power, Mint, nCr, sieve, factorization, Pollard rho, divisor sieve, totient, M\"obius,
  floor blocks, Pascal, extended gcd, CRT, matrix power, FFT, NTT, XOR basis, Grundy numbers, inclusion--exclusion \\
\file{string/} & hashing, KMP, Z-function, trie, binary trie, Aho--Corasick, suffix array + LCP, suffix automaton, Manacher \\
\file{geometry/} & point primitives, segment intersection, polygons (area, containment, Pick), convex hull, closest
  pair, rectangles \\
\file{flow/} & Dinic + min cut, min-cost flow, flows with lower bounds, Hopcroft--Karp \\
\file{misc/} & random numbers, interactive-problem helpers \\
\bottomrule
\end{tabularx}
\end{center}

\chapter{Maintenance and Troubleshooting}\label{ch:maintenance}

\section{VS Code tasks}

All tools are VS Code tasks (\key{Ctrl+Shift+P}, \emph{Tasks: Run Task}), defined in \file{.vscode/tasks.json}.
Each runs \file{tools/cp.py} with Python.

\begin{center}
\begin{tabularx}{\linewidth}{>{\raggedright\arraybackslash}p{4.2cm}p{3.3cm}L}
\toprule
\thd{Task} & \thd{Shortcut} & \thd{Command} \\
\midrule
\texttt{cf-run} & \key{F5} & \texttt{cp.py test FILE} \\
CP: Run With Typed Input & \key{Ctrl+F5} & \texttt{cp.py run FILE} \\
CP: Stress Test & \key{Ctrl+Alt+X} & \texttt{cp.py stress FILE} \\
CP: Quickstart & \key{Ctrl+Shift+N} & \texttt{cp.py new NAME} \\
CP: Quickstart USACO & \key{Ctrl+Shift+U} & \texttt{cp.py new NAME -{}-usaco} \\
CP: Archive Old Files & \key{Ctrl+Shift+Alt+C} & \texttt{cp.py archive} \\
CP: Setup & --- & \texttt{cp.py setup -{}-check} \\
CP: Build Docs & --- & \texttt{cp.py docs} \\
\bottomrule
\end{tabularx}
\end{center}

\section{CP: Setup}

Run it after changing anything in \file{library/}, the template, or the flags. It is idempotent, and does all of this:

\begin{enumerate}
  \item rebuilds the precompiled header for the current compiler and flags;
  \item regenerates \file{snippets/cp-library.code-snippets} in the \emph{competitive} profile;
  \item regenerates the shortcut block in the profile's \file{keybindings.json};
  \item writes the four CPH settings (flags, template path, template variables, stderr handling);
  \item regenerates \file{CHEATSHEET.md} and the generated parts of this guide (reference cards, keyboard maps);
  \item with \texttt{-{}-check}: compiles every template as strict C++17, both as written and under
        \cc{#define int long long}, and reports any warning.
\end{enumerate}

Reload VS Code once afterwards (\key{Ctrl+Shift+P}, \emph{Reload Window}) so new shortcuts take effect.

\section{Adding your own template}

Create a \file{.cpp} file in the right category folder. The header comments describe it; the code between
\texttt{// @begin} and \texttt{// @end} is what gets pasted. Anything outside the markers exists only so the file
compiles on its own (includes, a wrapper function providing the variables an \emph{inside solve()} snippet uses).

\begin{code}
// @name   Sum of a range (example)
// @prefix cprangesum
// @key    y z
// @desc   [global] range_sum(P, l, r) = sum of a[l, r) from prefix sums P.
// @usage  ll s = range_sum(P, l, r);
// @complexity O(1)
#include <bits/stdc++.h>
using namespace std;
using ll = long long;

// @begin
ll range_sum(const vector<ll>& P, int l, int r) { return P[r] - P[l]; }
// @end
\end{code}

\begin{itemize}
  \item \texttt{@key} is \texttt{t}, \texttt{y} or \texttt{u} (the leader) followed by a VS Code key name: a letter, a digit,
        \texttt{shift+x}, or one of \texttt{-~=~[~]~;~'~,~.~/~\textasciigrave}. Setup refuses duplicate keys, prefixes and names.
  \item \texttt{@desc} starts with \texttt{[global]} or \texttt{[inside solve]}; the rest is the one-line description.
  \item \texttt{@usage} and \texttt{@complexity} are optional and feed the reference cards.
  \item Do not end a \texttt{//} comment with a backslash (C++ would glue the next line into the comment).
\end{itemize}

Then run \emph{CP: Setup} (with its check) and reload VS Code.

\section{Rebuilding these PDFs}

\emph{CP: Build Docs} (\texttt{python tools/cp.py docs}) regenerates the reference and compiles
\file{docs/guide.tex} and \file{docs/algorithms.tex} with TinyTeX's \texttt{pdflatex} (three passes, for the table of
contents), writing \file{docs/GUIDE.pdf} and \file{docs/ALGORITHMS.pdf}. Build files go to \file{docs/build/}.
\texttt{-{}-only guide} or \texttt{-{}-only algorithms} builds one of them. The shared look lives in
\file{docs/cpstyle.sty}.

\section{Troubleshooting}

\begin{description}[style=nextline, font=\sffamily\bfseries\color{cpblue}]
  \item[A shortcut does nothing]
    Reload the window. Check you are in the \emph{competitive} profile (bottom-left gear) and in a \texttt{.cpp} file:
    the template shortcuts only work in C++ editors.
  \item[CPH says ``Compilation error'' about \texttt{debug.h}]
    CPH's flags are missing \texttt{-I.../include}. Run \emph{CP: Setup}.
  \item[CPH marks a correct test as failed]
    If its Standard Error box shows output, \texttt{ignoreSTDERROR} was turned off; run \emph{CP: Setup}. Otherwise the
    program exited with a non-zero code: press \key{F5} to see why.
  \item[Compiles are slow again]
    The compiler or flags changed and the precompiled header is stale. Run \emph{CP: Setup} (F5 rebuilds it
    automatically too).
  \item[GCC ``internal compiler error'' mentioning \texttt{bits/stdc++.h}]
    CPH's flags were edited by hand and no longer match the precompiled header. Run \emph{CP: Setup}.
  \item[Red squiggles under \texttt{bits/stdc++.h}]
    IntelliSense is configured in \file{.vscode/settings.json} to use MSYS2's \texttt{g++} in C++17 mode with
    \texttt{LOCAL} defined; if you moved the compiler, update \texttt{C\_Cpp.default.compilerPath}.
  \item[A deep recursion crashes locally]
    It should not with a 256\,MB stack; if it does, the recursion is deeper than about $10^6$ levels or each level holds a
    big local array. Make it iterative or move the array out of the function.
  \item[The stress tester says brute.cpp fails a sample]
    Fix the brute force first; a wrong brute force makes every verdict meaningless.
  \item[I changed the template and new files look old]
    CPH reads \file{template/template.cpp} at file creation time; existing files are never changed.
\end{description}

\section{Undoing everything}

\begin{enumerate}
  \item In the profile folder, restore \file{keybindings.json\allowbreak.bak-cp-setup} and
        \file{settings.json\allowbreak.bak-cp-setup} over the current files, and delete
        \file{snippets/\allowbreak cp-library.code-snippets}.
  \item Restore \file{template/template.cpp.bak} over \file{template/template.cpp}.
  \item Delete \file{.vscode/}, \file{tools/}, \file{include/}, \file{library/}, \file{stress/}, \file{docs/} and
        \file{CHEATSHEET.md} from the folder.
\end{enumerate}
Your solutions and CPH's test cases are never touched by any of the tools except the archive shortcut, which only
moves them.

\chapter{Library Reference}\label{ch:reference}
Every template in \file{library/}, generated from the files themselves by \cc{python tools/cp.py setup}, so this
chapter is always in sync with what the shortcuts paste. Each card shows the shortcut, the snippet prefix, where
the code belongs, its complexity and the source file, followed by a description and typical usage.
The algorithms behind each template are explained in the companion handbook, \emph{ALGORITHMS.pdf}.
% generated by tools/cp.py setup -- do not edit
\section{Template}
\begin{libentry}{Full template}{Ctrl+Alt+T, T}{cptemplate}{empty file}{--}{template/template.cpp}
the whole template, for files CPH didn\textquotesingle{}t create
\end{libentry}
\section{Basics: input, search, sorting, prefix sums, two pointers}
\begin{libentry}{2D prefix sums}{Ctrl+Alt+T, Shift+A}{cpprefix2d}{inside solve()}{O(nm) build, O(1) per query}{library/basics/prefix\_sum\_2d.cpp}
needs n, m, a (n x m). rect(r1, c1, r2, c2) = sum of rows [r1, r2) x cols [c1, c2).
\begin{lstlisting}
ll s = rect(r1, c1, r2, c2);  // rows [r1, r2), cols [c1, c2)
\end{lstlisting}
\end{libentry}
\begin{libentry}{All permutations (next\_permutation)}{Ctrl+Alt+T, 9}{cpperm}{inside solve()}{O(n! n)}{library/basics/permutations.cpp}
needs n (n \textless{}= \textasciitilde{}10). Visits every order of 0..n-1 in lexicographic order.
\begin{lstlisting}
// on a sorted vector v directly: do { ... } while (next_permutation(v.begin(), v.end()));  (skips duplicates)
\end{lstlisting}
\end{libentry}
\begin{libentry}{All subsets (bitmask loop)}{Ctrl+Alt+T, 8}{cpsubsets}{inside solve()}{O(2\textasciicircum{}n n)}{library/basics/subsets.cpp}
needs n, a (n \textless{}= \textasciitilde{}20). Visits every subset; bit i of mask set = item i taken.
\begin{lstlisting}
// for all submasks of a mask: for (int s = mask; s; s = (s - 1) & mask)
\end{lstlisting}
\end{libentry}
\begin{libentry}{Argsort + custom comparator}{Ctrl+Alt+T, Shift+J}{cpargsort}{above solve()}{O(n log n)}{library/basics/argsort.cpp}
argsort(a): indices sorted by a[i] (stable). Custom orders: sort with a lambda (see usage).
\begin{lstlisting}
vector<int> idx = argsort(a);  // a[idx[0]] is the smallest
sort(v.begin(), v.end(), [](const auto& x, const auto& y) { return make_pair(x.first, -x.second) < make_pair(y.first, -y.second); });
\end{lstlisting}
\end{libentry}
\begin{libentry}{Binary search on a real answer}{Ctrl+Alt+T, Shift+F}{cpbsreal}{inside solve()}{100 calls to ok}{library/basics/binary\_search\_real.cpp}
smallest real x in [lo, hi] with ok(x); ok must be monotone. 100 halvings reach full double precision.
\begin{lstlisting}
// fill in ok(x); lo ends within 1e-15 of the boundary (relative to hi - lo)
\end{lstlisting}
\end{libentry}
\begin{libentry}{Binary search on the answer}{Ctrl+Alt+T, A}{cpbsearch}{inside solve()}{O(log(hi - lo)) calls to ok}{library/basics/binary\_search.cpp}
smallest x in [lo, hi] with ok(x); ok must be monotone (false...false true...true).
\begin{lstlisting}
// fill in ok(x); lo ends as the smallest x with ok(x)
\end{lstlisting}
\end{libentry}
\begin{libentry}{Bit helpers}{Ctrl+Alt+T, Y}{cpbits}{above solve()}{O(1)}{library/basics/bits.cpp}
popcnt, highest\_bit (floor log2), lowest\_bit, has\_bit for 64-bit values; submask loop in usage.
\begin{lstlisting}
popcnt(x); highest_bit(x); lowest_bit(x); has_bit(x, i);
for (ll s = mask; s; s = (s - 1) & mask) { /* every non-empty submask */ }
\end{lstlisting}
\end{libentry}
\begin{libentry}{Coordinate compression}{Ctrl+Alt+T, Shift+C}{cpcompress}{above solve()}{O(n log n)}{library/basics/compress.cpp}
auto vals = compress(a); a[i] becomes its rank 0..k-1 and vals[rank] is the original value.
\begin{lstlisting}
vector<int> vals = compress(a);  // a[i] -> rank, vals[rank] -> value
\end{lstlisting}
\end{libentry}
\begin{libentry}{Difference array (offline range add, 1D)}{Ctrl+Alt+T, 3}{cpdiff}{inside solve()}{O(1) per update, O(n) build}{library/basics/difference\_array.cpp}
needs n. range\_add(l, r, v) adds v to a[l, r) in O(1); build() returns the final array.
\begin{lstlisting}
range_add(l, r, v);            // any number of times
vector<ll> a = build();        // once, at the end
\end{lstlisting}
\end{libentry}
\begin{libentry}{Difference array (offline rectangle add, 2D)}{Ctrl+Alt+T, 4}{cpdiff2d}{inside solve()}{O(1) per update, O(nm) build}{library/basics/difference\_array\_2d.cpp}
needs n, m. rect\_add(r1, c1, r2, c2, v) adds v to rows [r1, r2) x cols [c1, c2); build() returns the grid.
\begin{lstlisting}
rect_add(r1, c1, r2, c2, v);   // any number of times
auto g = build();              // g[r][c] = total added to cell (r, c)
\end{lstlisting}
\end{libentry}
\begin{libentry}{Gray code}{Ctrl+Alt+U, 0}{cpgray}{above solve()}{O(1) / O(log) per value}{library/basics/gray\_code.cpp}
gray(i) = i \textasciicircum{} (i \textgreater{}\textgreater{} 1): consecutive values differ in exactly one bit; gray\_inverse(g) gives i back. To visit all 2\textasciicircum{}n subsets changing one element at a time, step k flips bit \_\_builtin\_ctzll(k).
\begin{lstlisting}
for (ll k = 1, mask = 0; k < (1LL << n); k++) { int b = __builtin_ctzll(k); mask ^= 1LL << b; /* add or remove b */ }
ll g = gray(i), back = gray_inverse(g);
\end{lstlisting}
\end{libentry}
\begin{libentry}{Intervals: merge + max disjoint (greedy)}{Ctrl+Alt+T, Shift+N}{cpintervals}{above solve()}{O(n log n)}{library/basics/intervals.cpp}
merge\_intervals(v): union of closed [l, r] intervals; max\_disjoint(v): most pairwise non-overlapping (earliest end first).
\begin{lstlisting}
vector<pair<ll, ll>> segs = merge_intervals(v);
int k = max_disjoint(v);  // touching intervals ([1, 2] and [2, 3]) count as overlapping
\end{lstlisting}
\end{libentry}
\begin{libentry}{Meet in the middle (subset sums)}{Ctrl+Alt+T, Shift+U}{cpmitm}{above solve()}{O(2\textasciicircum{}(n/2) n)}{library/basics/meet\_in\_middle.cpp}
count\_subsets\_leq(a, X): number of subsets with sum \textless{}= X for n \textless{}= \textasciitilde{}40, splitting the items in two halves.
\begin{lstlisting}
ll ways = count_subsets_leq(a, X);
// exact sum X: count upper_bound - lower_bound of X - x in the sorted right half
\end{lstlisting}
\end{libentry}
\begin{libentry}{Monotonic stack (next greater / previous smaller)}{Ctrl+Alt+T, 7}{cpmonostack}{above solve()}{O(n)}{library/basics/monotonic\_stack.cpp}
next\_greater(a)[i] = first j \textgreater{} i with a[j] \textgreater{} a[i] (n if none); prev\_smaller(a)[i] = last j \textless{} i with a[j] \textless{} a[i] (-1).
\begin{lstlisting}
auto nxt = next_greater(a); auto prv = prev_smaller(a);
// largest rectangle in a histogram: width of bar i = next smaller - prev smaller - 1
\end{lstlisting}
\end{libentry}
\begin{libentry}{Prefix sums (1D)}{Ctrl+Alt+T, 1}{cpprefix}{inside solve()}{O(n) build, O(1) per query}{library/basics/prefix\_sum.cpp}
needs n, a. P[i] = a[0] + ... + a[i-1]; range\_sum(l, r) = a[l] + ... + a[r-1].
\begin{lstlisting}
ll s = range_sum(l, r);  // [l, r)
// subarrays with sum k: count pairs i < j with P[j] - P[i] == k using a map of seen P values
\end{lstlisting}
\end{libentry}
\begin{libentry}{Radix sort and counting sort}{Ctrl+Alt+U, Shift+R}{cpradix}{above solve()}{O(n) per pass: O(4 (n + 65536)) for radix\_sort, O(n + K) for counting\_order}{library/basics/radix\_sort.cpp}
radix\_sort(a): sorts 64-bit integers (negatives too) in 4 linear passes of 16 bits, often 2-3x faster than std::sort for n \textgreater{}= 1e6; counting\_order(keys, K): stable order of indices by small keys in [0, K).
\begin{lstlisting}
radix_sort(a);   // vector<ll>
vector<int> ord = counting_order(keys, K);   // stable: equal keys keep their order
\end{lstlisting}
\end{libentry}
\begin{libentry}{Read graph (n, m, m edges)}{Ctrl+Alt+T, Shift+G}{cpgraph}{inside solve()}{O(n + m)}{library/basics/read\_graph.cpp}
reads n, m and m 1-indexed edges into a 0-indexed undirected adjacency list.
\begin{lstlisting}
// then: dfs / bfs / ... over adj
\end{lstlisting}
\end{libentry}
\begin{libentry}{Read vector}{Ctrl+Alt+T, V}{cpread}{inside solve()}{O(n)}{library/basics/read\_vector.cpp}
needs n. Reads n numbers into a.
\begin{lstlisting}
// with the template: vi a(n); cin >> a;  does the same
\end{lstlisting}
\end{libentry}
\begin{libentry}{Read weighted graph (n, m, edges u v w)}{Ctrl+Alt+T, Shift+E}{cpgraphw}{inside solve()}{O(n + m)}{library/basics/read\_weighted\_graph.cpp}
reads n, m and m 1-indexed weighted edges into adj[u] = \{v, w\}, 0-indexed, undirected.
\begin{lstlisting}
// then: auto [dist, par] = dijkstra(adj, 0);
\end{lstlisting}
\end{libentry}
\begin{libentry}{Recursive complete search (backtracking)}{Ctrl+Alt+T, Shift+Y}{cpbacktrack}{inside solve()}{O(2\textasciicircum{}n) calls without pruning}{library/basics/backtracking.cpp}
needs n. Builds every subsequence of 0..n-1 in cur (skip / take); prune at the top of search().
\begin{lstlisting}
// permutations instead: loop over unused items, mark used, recurse, unmark
\end{lstlisting}
\end{libentry}
\begin{libentry}{Sliding window min / max (monotonic deque)}{Ctrl+Alt+T, 6}{cpwindow}{above solve()}{O(n)}{library/basics/sliding\_window.cpp}
window\_min(a, k)[i] = min of a[i, i + k); pass want\_max = true for the max. Monotonic deque.
\begin{lstlisting}
vector<ll> mins = window_min(a, k);
vector<ll> maxs = window_min(a, k, true);
\end{lstlisting}
\end{libentry}
\begin{libentry}{Ternary search (unimodal max)}{Ctrl+Alt+T, Q}{cpternary}{above solve()}{O(log range) calls to f}{library/basics/ternary\_search.cpp}
argmax\_int(lo, hi, f) / argmax\_real(lo, hi, f) for f that rises then falls. For a minimum pass -f.
\begin{lstlisting}
ll x = argmax_int(0, 1e9, [&](ll x) { return profit(x); });
ld y = argmax_real(0, 1e4, [&](ld x) { return -cost(x); });  // minimum of cost
\end{lstlisting}
\end{libentry}
\begin{libentry}{Two pointers (sliding window)}{Ctrl+Alt+T, 5}{cptwoptr}{inside solve()}{O(n)}{library/basics/two\_pointers.cpp}
needs n, a (values \textgreater{}= 0). Longest window with sum \textless{}= S; grow r, shrink l while the window is invalid.
\begin{lstlisting}
// count windows instead: add (r - l + 1) to an answer inside the loop
\end{lstlisting}
\end{libentry}
\section{Graphs}
\begin{libentry}{0-1 BFS}{Ctrl+Alt+T, 0}{cp01bfs}{above solve()}{O(n + m)}{library/graph/zero\_one\_bfs.cpp}
vector\textless{}ll\textgreater{} dist = zero\_one\_bfs(adj, src); adj[u] = \{v, w\} with w in \{0, 1\}. Unreachable = LLONG\_MAX.
\begin{lstlisting}
vector<ll> dist = zero_one_bfs(adj, src);  // adj[u] = {v, w}, w in {0, 1}
\end{lstlisting}
\end{libentry}
\begin{libentry}{2-edge-connected components}{Ctrl+Alt+Y, B}{cp2ecc}{above solve()}{O(n + m)}{library/graph/two\_edge\_cc.cpp}
auto [k, comp] = two\_edge\_cc(n, edges): vertices stay connected after deleting any single edge iff same comp. Bridges join comps.
\begin{lstlisting}
auto [k, comp] = two_edge_cc(n, edges);
// bridge tree: for each edge (u, v) with comp[u] != comp[v], connect comp[u] - comp[v]
\end{lstlisting}
\end{libentry}
\begin{libentry}{2-SAT}{Ctrl+Alt+Y, 2}{cp2sat}{above solve()}{O(n + clauses)}{library/graph/two\_sat.cpp}
TwoSat ts(n); ts.either(a, va, b, vb) adds the clause (x\_a == va OR x\_b == vb); ts.solve() -\textgreater{} values or empty if impossible.
\begin{lstlisting}
TwoSat ts(n);
ts.either(a, true, b, false); ts.implies(a, true, c, true); ts.set(d, false);
vector<int> x = ts.solve();  // empty: unsatisfiable
\end{lstlisting}
\end{libentry}
\begin{libentry}{Bellman-Ford (negative edges)}{Ctrl+Alt+T, Shift+W}{cpbellman}{above solve()}{O(nm)}{library/graph/bellman\_ford.cpp}
bellman\_ford(n, edges, src), edges = \{u, v, w\}. BF\_INF = unreachable, -BF\_INF = hit by a negative cycle.
\begin{lstlisting}
vector<array<ll, 3>> edges;  // {u, v, w}
vector<ll> dist = bellman_ford(n, edges, src);  // -BF_INF: negative cycle reaches it
\end{lstlisting}
\end{libentry}
\begin{libentry}{BFS (shortest path, unweighted)}{Ctrl+Alt+T, B}{cpbfs}{inside solve()}{O(n + m)}{library/graph/bfs.cpp}
needs n, adj, src. dist[v] = -1 if unreachable; get\_path(t) rebuilds the route.
\begin{lstlisting}
// after it: dist[v] (-1 = unreachable), par[v]
vector<int> path = get_path(t);  // src ... t
\end{lstlisting}
\end{libentry}
\begin{libentry}{Biconnected components (blocks)}{Ctrl+Alt+Y, Shift+B}{cpbcc}{above solve()}{O(n + m)}{library/graph/biconnected\_components.cpp}
biconnected\_components(n, edges): every block (maximal subgraph with no cut vertex) as a vertex list. Cut vertices appear in several.
\begin{lstlisting}
auto blocks = biconnected_components(n, edges);
// block-cut tree: node per block + node per cut vertex, edge when the cut vertex is in the block
\end{lstlisting}
\end{libentry}
\begin{libentry}{Bidirectional BFS (meet in the middle on shortest paths)}{Ctrl+Alt+U, 2}{cpbibfs}{above solve()}{O(b\textasciicircum{}(d/2)) states for branching factor b and distance d, instead of O(b\textasciicircum{}d)}{library/graph/bidirectional\_bfs.cpp}
bidirectional\_bfs(s, t, neighbours): unweighted distance from s to t (-1 if unreachable) exploring from both ends; neighbours(state) returns a vector of states. Works on implicit graphs (puzzles, strings) with any hashable State.
\begin{lstlisting}
ll d = bidirectional_bfs<ll>(start, goal, [&](ll s) { vector<ll> nb; /* moves */ return nb; });
\end{lstlisting}
\end{libentry}
\begin{libentry}{Bipartite check (2-coloring)}{Ctrl+Alt+T, 2}{cpbipartite}{above solve()}{O(n + m)}{library/graph/bipartite.cpp}
vector\textless{}int\textgreater{} color = bipartite(adj); colors are 0/1, empty if an odd cycle exists.
\begin{lstlisting}
vector<int> color = bipartite(adj);
if (color.empty()) { /* odd cycle: not bipartite */ }
\end{lstlisting}
\end{libentry}
\begin{libentry}{Bridges + articulation points}{Ctrl+Alt+T, Shift+B}{cpbridges}{above solve()}{O(n + m)}{library/graph/bridges.cpp}
auto [bridge\_ids, cut\_vertices] = bridges\_and\_cuts(n, edges); undirected, multi-edges OK.
\begin{lstlisting}
vector<pair<int, int>> edges;  // undirected {u, v}
auto [bridge_ids, cut_vertices] = bridges_and_cuts(n, edges);
\end{lstlisting}
\end{libentry}
\begin{libentry}{Cycle finding in O(1) memory (Floyd / Brent)}{Ctrl+Alt+U, F}{cpcyclefind}{above solve()}{O(mu + lambda) evaluations of f, O(1) memory}{library/graph/cycle\_floyd.cpp}
find\_rho(x0, f): for the sequence x0, f(x0), f(f(x0)), ... returns \{mu, lambda\}: the first index of the cycle and its length. Brent\textquotesingle{}s variant: O(mu + lambda) calls of f, O(1) memory (use when states are huge).
\begin{lstlisting}
auto [mu, lam] = find_rho<ll>(seed, [](ll x) { return (x * x + 1) % 1'000'003; });
\end{lstlisting}
\end{libentry}
\begin{libentry}{DFS (recursive, all components)}{Ctrl+Alt+T, D}{cpdfs}{inside solve()}{O(n + m)}{library/graph/dfs.cpp}
needs n, adj. Labels every vertex with its component id.
\begin{lstlisting}
// after it: comp[v] = component of v, ncomp = number of components
\end{lstlisting}
\end{libentry}
\begin{libentry}{Dijkstra}{Ctrl+Alt+T, J}{cpdijkstra}{above solve()}{O((n + m) log n)}{library/graph/dijkstra.cpp}
auto [dist, par] = dijkstra(adj, src); adj[u] = \{v, w\}, w \textgreater{}= 0. Unreachable = LLONG\_MAX.
\begin{lstlisting}
vector<vector<pair<int, ll>>> adj(n);  // adj[u].push_back({v, w})
auto [dist, par] = dijkstra(adj, src);
\end{lstlisting}
\end{libentry}
\begin{libentry}{DSU (union-find)}{Ctrl+Alt+T, U}{cpdsu}{above solve()}{\textasciitilde{}O(1) amortized per call}{library/graph/dsu.cpp}
DSU d(n); d.unite(a, b) -\textgreater{} false if already joined; d.same, d.size, d.comps.
\begin{lstlisting}
DSU d(n);
if (d.unite(a, b)) { /* merged two components */ }
d.same(a, b); d.size(a); d.comps;
\end{lstlisting}
\end{libentry}
\begin{libentry}{DSU with rollback}{Ctrl+Alt+Y, R}{cpdsurb}{above solve()}{O(log n) per find / unite, O(1) per undone unite}{library/graph/dsu\_rollback.cpp}
DSURollback d(n); d.unite(a, b); int t = d.snapshot(); ... d.rollback(t) undoes every unite after t. No path compression.
\begin{lstlisting}
DSURollback d(n);
int t = d.snapshot(); d.unite(a, b); /* ... */ d.rollback(t);
\end{lstlisting}
\end{libentry}
\begin{libentry}{Eulerian path / circuit (Hierholzer)}{Ctrl+Alt+Y, E}{cpeuler}{above solve()}{O(n + m)}{library/graph/euler\_path.cpp}
euler\_path(n, edges, directed, src = -1): vertex sequence using every edge exactly once, or empty if impossible.
\begin{lstlisting}
vector<int> path = euler_path(n, edges, false);       // path.size() == m + 1
vector<int> tour = euler_path(n, edges, false, 0);    // circuit from 0 (only when every degree is even)
\end{lstlisting}
\end{libentry}
\begin{libentry}{Find a directed cycle}{Ctrl+Alt+T, /}{cpcycle}{above solve()}{O(n + m)}{library/graph/find\_cycle.cpp}
find\_cycle(adj): vertices of some directed cycle in order (the last one has an edge back to the first); empty if acyclic.
\begin{lstlisting}
vector<int> cyc = find_cycle(adj);
// undirected graphs: a DSU finds the first edge whose ends are already joined
\end{lstlisting}
\end{libentry}
\begin{libentry}{Flood fill (grid regions)}{Ctrl+Alt+T, F}{cpflood}{inside solve()}{O(nm)}{library/graph/flood\_fill.cpp}
needs n, m, grid (vector\textless{}string\textgreater{}). Iterative, so big grids can\textquotesingle{}t overflow the stack.
\begin{lstlisting}
// after it: region[r][c] = region id (-1 = not open), region_size[id]
// change open(r, c) to choose which cells get filled
\end{lstlisting}
\end{libentry}
\begin{libentry}{Floyd-Warshall (all pairs)}{Ctrl+Alt+T, W}{cpfloyd}{above solve()}{O(n\textasciicircum{}3)}{library/graph/floyd\_warshall.cpp}
floyd\_warshall(d) in place, O(n\textasciicircum{}3). d[i][i] \textless{} 0 afterwards \textless{}=\textgreater{} i is on a negative cycle.
\begin{lstlisting}
vector<vector<ll>> d(n, vector<ll>(n, FW_INF));
for (int i = 0; i < n; i++) d[i][i] = 0;  // then d[u][v] = min(d[u][v], w) per edge
floyd_warshall(d);
\end{lstlisting}
\end{libentry}
\begin{libentry}{Functional graph (successor graph) cycles}{Ctrl+Alt+T, ,}{cpfuncgraph}{above solve()}{O(n)}{library/graph/functional\_graph.cpp}
FunctionalGraph fg(f), f[v] = the one successor of v: fg.cyc[v] (cycle id or -1), fg.cyc\_len[id], fg.dist[v] to its cycle, fg.entry[v].
\begin{lstlisting}
FunctionalGraph fg(f);
int steps_until_repeat = fg.dist[v] + fg.cyc_len[fg.cyc[fg.entry[v]]];
\end{lstlisting}
\end{libentry}
\begin{libentry}{Grid BFS (shortest path on a grid)}{Ctrl+Alt+T, G}{cpgridbfs}{inside solve()}{O(nm)}{library/graph/grid\_bfs.cpp}
needs n, m, grid. Starts from every \textquotesingle{}S\textquotesingle{} (multi-source), \textquotesingle{}\#\textquotesingle{} is a wall.
\begin{lstlisting}
// after it: dist[r][c] = steps from the nearest 'S' (-1 = unreachable)
\end{lstlisting}
\end{libentry}
\begin{libentry}{K-th successor (binary lifting / jumping)}{Ctrl+Alt+T, .}{cpjump}{above solve()}{O(n log K) build, O(log K) per query}{library/graph/kth\_successor.cpp}
Jump J(f, max\_k); J.go(v, k) = f applied k times to v, for k \textless{}= max\_k. Works for any functional graph.
\begin{lstlisting}
Jump J(f, 1e9);
int where = J.go(v, k);
\end{lstlisting}
\end{libentry}
\begin{libentry}{Kruskal MST}{Ctrl+Alt+T, K}{cpkruskal}{above solve()}{O(m log m)}{library/graph/kruskal.cpp}
auto [cost, used] = kruskal(n, edges); edges = \{u, v, w\}. Spanning forest if disconnected.
\begin{lstlisting}
vector<array<ll, 3>> edges;  // {u, v, w}
auto [cost, used] = kruskal(n, edges);
\end{lstlisting}
\end{libentry}
\begin{libentry}{Offline dynamic connectivity (edge deletions)}{Ctrl+Alt+Y, Shift+R}{cpofflinecc}{above solve()}{O(q log q log n)}{library/graph/offline\_connectivity.cpp}
offline\_components(n, ops): ops are \{0, u, v\} add edge, \{1, u, v\} remove it, \{2, 0, 0\} query; returns \#components per query. Needs cpdsurb.
\begin{lstlisting}
vector<array<int, 3>> ops;  // {type, u, v}
vector<int> answers = offline_components(n, ops);
\end{lstlisting}
\end{libentry}
\begin{libentry}{Prim MST (dense graphs)}{Ctrl+Alt+T, Shift+5}{cpprim}{above solve()}{O(n\textasciicircum{}2)}{library/graph/prim.cpp}
prim\_dense(w): MST weight for an n x n weight matrix (complete graphs, e.g. points in the plane). O(n\textasciicircum{}2).
\begin{lstlisting}
vector<vector<ll>> w(n, vector<ll>(n));  // w[i][j] = cost of edge i-j
ll cost = prim_dense(w);
\end{lstlisting}
\end{libentry}
\begin{libentry}{SCC (Tarjan)}{Ctrl+Alt+T, C}{cpscc}{above solve()}{O(n + m)}{library/graph/scc.cpp}
auto [k, comp] = scc(adj). Ids are reverse-topological: edge u-\textgreater{}v across SCCs has comp[u] \textgreater{} comp[v].
\begin{lstlisting}
auto [k, comp] = scc(adj);  // k components, comp[v] in [0, k)
\end{lstlisting}
\end{libentry}
\begin{libentry}{System of difference constraints}{Ctrl+Alt+U, Shift+I}{cpdiffcons}{above solve()}{O(n m) (Bellman-Ford from a virtual source)}{library/graph/difference\_constraints.cpp}
DiffConstraints dc(n); dc.add(u, v, w) means x\_v - x\_u \textless{}= w; dc.solve() -\textgreater{} a solution vector (all x \textless{}= 0, maximal), or empty if infeasible (a negative cycle).
\begin{lstlisting}
DiffConstraints dc(n); dc.add(a, b, 3);   // x_b <= x_a + 3
vector<ll> x = dc.solve(); if (x.empty()) { /* impossible */ }
\end{lstlisting}
\end{libentry}
\begin{libentry}{Topological sort (Kahn)}{Ctrl+Alt+T, O}{cptopo}{above solve()}{O(n + m)}{library/graph/topo\_sort.cpp}
vector\textless{}int\textgreater{} order = topo\_sort(adj); empty if the graph has a cycle.
\begin{lstlisting}
vector<int> order = topo_sort(adj);
if (order.empty()) { /* the graph has a cycle */ }
\end{lstlisting}
\end{libentry}
\section{Trees}
\begin{libentry}{Centroid decomposition}{Ctrl+Alt+Y, C}{cpcentroid}{above solve()}{O(n log n)}{library/tree/centroid.cpp}
CentroidDecomp cd(adj): cd.cpar[v] = parent in the centroid tree (-1 at the top), cd.clevel[v] = its depth (\textless{}= log2 n).
\begin{lstlisting}
CentroidDecomp cd(adj);
for (int c = v; c != -1; c = cd.cpar[c]) { /* every path through v passes a centroid ancestor c: use dist(v, c) */ }
\end{lstlisting}
\end{libentry}
\begin{libentry}{DSU on tree (Sack): subtree statistics in one shared array}{Ctrl+Alt+U, Shift+T}{cpsack}{inside solve()}{O(n log n) calls of add}{library/tree/sack.cpp}
needs n, adj, root, color (values in [0, n)). Keeps the heavy child\textquotesingle{}s counts and re-adds the light subtrees: best[v] = the largest count of one color in v\textquotesingle{}s subtree. Edit add / clear / the answer line.
\begin{lstlisting}
// add(v) puts one vertex into cnt[] and the statistic; clear(v) empties the subtree of v again
\end{lstlisting}
\end{libentry}
\begin{libentry}{Heavy-light decomposition}{Ctrl+Alt+Y, H}{cphld}{above solve()}{O(n) build, O(log n) ranges per path}{library/tree/hld.cpp}
HLD h(adj, root); h.for\_path(a, b, f) calls f(l, r) on position ranges [l, r) covering the path; pair with any [l, r) segment tree.
\begin{lstlisting}
HLD h(adj); SegTree st(n);                      // st.set(h.pos[v], value of v)
ll s = 0; h.for_path(a, b, [&](int l, int r) { s += st.query(l, r); });
// subtree of v = [h.pos[v], h.pos[v] + h.sub[v]); h.lca(a, b)
\end{lstlisting}
\end{libentry}
\begin{libentry}{Jump pointers (level ancestors with one extra pointer per vertex)}{Ctrl+Alt+U, J}{cpjumpptr}{above solve()}{O(1) per added leaf, O(log n) per query}{library/tree/jump\_pointers.cpp}
JumpPointers jp(n); jp.add\_root(r); jp.add\_leaf(v, p) as vertices arrive (parent first); jp.kth(v, k), jp.ancestor\_at\_depth(v, d), jp.lca(a, b). Works online while the tree grows.
\begin{lstlisting}
JumpPointers jp(n); jp.add_root(0);
for (int v : order) if (v != 0) jp.add_leaf(v, par[v]);  // any order with parents before children
int a = jp.kth(v, k);  // -1 if v has fewer than k ancestors
\end{lstlisting}
\end{libentry}
\begin{libentry}{Kruskal reconstruction tree}{Ctrl+Alt+Y, K}{cpkrt}{above solve()}{O(m log m)}{library/tree/kruskal\_reconstruction.cpp}
KRT t(n, edges \{u, v, w\}): leaves 0..n-1, each merge makes node x with val[x] = w. Max edge on the MST path a-b = val[lca(a, b)].
\begin{lstlisting}
KRT t(n, edges); auto g = t.tree();  // then LCA l(g, g.size() - 1) if the graph is connected
// vertices reachable from v using edges <= x = leaves under v's highest ancestor with val <= x
\end{lstlisting}
\end{libentry}
\begin{libentry}{LCA (binary lifting)}{Ctrl+Alt+T, L}{cplca}{above solve()}{O(n log n) build, O(log n) per query}{library/tree/lca.cpp}
LCA t(adj, root); t.lca(a, b), t.dist(a, b), t.kth\_ancestor(v, k), t.depth[v].
\begin{lstlisting}
LCA t(adj, root);
t.lca(a, b); t.dist(a, b); t.kth_ancestor(v, k); t.depth[v];
\end{lstlisting}
\end{libentry}
\begin{libentry}{Long-path decomposition: depth-indexed tree DP in O(n)}{Ctrl+Alt+U, Shift+N}{cplongpath}{inside solve()}{O(n)}{library/tree/long\_path.cpp}
needs n, adj, root. f\_v[d] = \# vertices at distance d below v lives at buf[pos[v] + d], shared with the long child (shifted by one); dom[v] = the smallest d maximising f\_v[d]. Edit the merge for other depth DPs.
\begin{lstlisting}
// the array of v has h[v] cells: buf[pos[v]] .. buf[pos[v] + h[v] - 1]
\end{lstlisting}
\end{libentry}
\begin{libentry}{Offline LCA (Tarjan)}{Ctrl+Alt+U, Shift+L}{cpofflinelca}{above solve()}{O((n + q) alpha(n))}{library/tree/offline\_lca.cpp}
tarjan\_lca(adj, root, qs): the LCA of every query pair, answered in one DFS with a DSU. Faster than binary lifting when all queries are known in advance and memory is tight.
\begin{lstlisting}
vector<int> ans = tarjan_lca(adj, 0, qs);   // qs: vector<pair<int, int>>
\end{lstlisting}
\end{libentry}
\begin{libentry}{Prufer code (tree \textless{}-\textgreater{} sequence)}{Ctrl+Alt+U, U}{cpprufer}{above solve()}{O(n) (linear pointer technique)}{library/tree/prufer.cpp}
prufer\_encode(n, edges): the Prufer sequence (length n-2) of a labelled tree on 0..n-1; prufer\_decode(seq): the tree back. Bijection behind Cayley\textquotesingle{}s formula n\textasciicircum{}(n-2); vertex v appears deg(v)-1 times.
\begin{lstlisting}
vector<int> seq = prufer_encode(n, edges);
vector<pair<int, int>> tree = prufer_decode(seq);
\end{lstlisting}
\end{libentry}
\begin{libentry}{Rerooting (DP for every root)}{Ctrl+Alt+T, \textquotesingle{}}{cpreroot}{above solve()}{O(n) merges and lifts}{library/tree/rerooting.cpp}
reroot\textless{}T\textgreater{}(adj, id, merge, lift): res[v] = merge over v\textquotesingle{}s neighbours w of lift(value of w\textquotesingle{}s side, w). merge associative, id its identity.
\begin{lstlisting}
// farthest vertex from every v:  auto far = reroot<ll>(adj, 0, [](ll a, ll b) { return max(a, b); }, [](ll x, int) { return x + 1; });
// sum of distances: T = pair<ll, ll> {sum, count}, lift(p, w) = {p.first + p.second + 1, p.second + 1}, merge = add
\end{lstlisting}
\end{libentry}
\begin{libentry}{Small-to-large merging}{Ctrl+Alt+Y, M}{cpsmall2large}{inside solve()}{O(n log\textasciicircum{}2 n)}{library/tree/small\_to\_large.cpp}
needs n, adj, root, color. distinct[v] = \# distinct colors in v\textquotesingle{}s subtree; always pour the smaller set into the larger.
\begin{lstlisting}
// swap set<int> for map<int, int> to keep counts (e.g. most frequent color per subtree)
\end{lstlisting}
\end{libentry}
\begin{libentry}{Tree DFS (parent, depth, subtree size, Euler tour)}{Ctrl+Alt+T, E}{cptree}{inside solve()}{O(n)}{library/tree/tree\_dfs.cpp}
needs n, adj, root. Subtree of u = Euler positions [tin[u], tout[u]).
\begin{lstlisting}
// after it: par, depth, sub (subtree size), tin/tout (Euler tour)
// subtree of u = positions [tin[u], tout[u]) -> a range query on the tour
\end{lstlisting}
\end{libentry}
\begin{libentry}{Tree diameter}{Ctrl+Alt+T, Shift+D}{cpdiameter}{above solve()}{O(n)}{library/tree/tree\_diameter.cpp}
auto [len, path] = tree\_diameter(adj); two BFS passes, path lists the vertices end to end.
\begin{lstlisting}
auto [len, path] = tree_diameter(adj);
\end{lstlisting}
\end{libentry}
\begin{libentry}{Tree DP (subtree DP, iterative)}{Ctrl+Alt+T, ;}{cptreedp}{inside solve()}{O(n)}{library/tree/tree\_dp.cpp}
needs n, adj, root. Children before parents (reverse BFS order); example: maximum independent set.
\begin{lstlisting}
// answer: max(dp[root][0], dp[root][1]); change the two lines marked "transition"
\end{lstlisting}
\end{libentry}
\begin{libentry}{Tree isomorphism (canonical IDs, AHU)}{Ctrl+Alt+U, H}{cptreehash}{above solve()}{O(n log n)}{library/tree/tree\_hash.cpp}
TreeCanon tc; tc.rooted(adj, root) and tc.unrooted(adj) give integer IDs: two trees get the same ID (from the same TreeCanon object) iff they are isomorphic. rooted\_all(adj, root)[v] is the ID of v\textquotesingle{}s subtree.
\begin{lstlisting}
TreeCanon tc;
bool iso = tc.unrooted(adj1) == tc.unrooted(adj2);
\end{lstlisting}
\end{libentry}
\begin{libentry}{Tree knapsack (merging child DPs)}{Ctrl+Alt+Y, Shift+K}{cptreeknap}{inside solve()}{O(n\textasciicircum{}2)}{library/tree/tree\_knapsack.cpp}
needs n, adj, root, w. best[v][j] = max value choosing j vertices of v\textquotesingle{}s subtree, each chosen vertex\textquotesingle{}s parent chosen too.
\begin{lstlisting}
// answer for exactly k vertices containing the root: best[root][k]; merging costs O(n^2) in total
\end{lstlisting}
\end{libentry}
\begin{libentry}{Virtual (auxiliary) tree}{Ctrl+Alt+Y, V}{cpvirtual}{above solve()}{O(k log k) LCA calls}{library/tree/virtual\_tree.cpp}
virtual\_tree(nodes, tin, lca) -\textgreater{} \{vertices, parent\}: key nodes + pairwise LCAs, parent[i] = virtual parent of vertices[i] (-1 root).
\begin{lstlisting}
auto [vs, vpar] = virtual_tree(keys, tin, [&](int a, int b) { return t.lca(a, b); });  // tin from cptree, t from cplca
// then DP over vs in reverse order: children appear after their parents
\end{lstlisting}
\end{libentry}
\section{Data structures \& range queries}
\begin{libentry}{Anti-hack hash for unordered\_map}{Ctrl+Alt+T, Shift+H}{cphash}{above solve()}{O(1) expected per operation}{library/ds/custom\_hash.cpp}
unordered\_map\textless{}ll, int, custom\_hash\textgreater{} mp; also works for pair keys. Stops anti-hash TLE hacks.
\begin{lstlisting}
unordered_map<ll, int, custom_hash> cnt;
unordered_map<pair<int, int>, int, custom_hash> seen;
\end{lstlisting}
\end{libentry}
\begin{libentry}{Cartesian tree}{Ctrl+Alt+U, C}{cpcartesian}{above solve()}{O(n) (one stack pass)}{library/ds/cartesian\_tree.cpp}
CartesianTree ct(a): min-heap on values, in-order = array order (ties: the leftmost is higher). ct.root, ct.par[i], ct.lc[i], ct.rc[i] (-1 = none). The subtree of i is the maximal range where a[i] is the minimum.
\begin{lstlisting}
CartesianTree ct(a);
int r = ct.root; for (int i = 0; i < n; i++) { int p = ct.par[i]; }
\end{lstlisting}
\end{libentry}
\begin{libentry}{CDQ divide and conquer (3D dominance counting)}{Ctrl+Alt+U, 3}{cpcdq}{above solve()}{O(n log\textasciicircum{}2 n)}{library/ds/cdq.cpp}
cdq\_dominance(pts): for each point i, the number of j != i with a\_j \textless{}= a\_i, b\_j \textless{}= b\_i, c\_j \textless{}= c\_i. The pattern (sort by one key, divide and conquer on the second, BIT on the third) solves offline 3D problems.
\begin{lstlisting}
vector<array<ll, 3>> pts(n);   // (a, b, c)
vector<int> dom = cdq_dominance(pts);
\end{lstlisting}
\end{libentry}
\begin{libentry}{Disjoint sparse table (O(1) queries for any associative operation)}{Ctrl+Alt+U, T}{cpdisjointsparse}{above solve()}{O(n log n) build, O(1) query}{library/ds/disjoint\_sparse\_table.cpp}
DisjointSparse\textless{}T, Op\textgreater{} t(a, op); t.query(l, r) = a[l] op ... op a[r-1] for [l, r), l \textless{} r. Works for sums, products mod p, matrix products: no need for idempotence.
\begin{lstlisting}
DisjointSparse t(a, [](ll x, ll y) { return x * y % MOD; });
ll prod = t.query(l, r);
\end{lstlisting}
\end{libentry}
\begin{libentry}{Distinct values in ranges (offline sweep + BIT)}{Ctrl+Alt+Y, Shift+D}{cpdistinct}{above solve()}{O((n + q) log n)}{library/ds/distinct\_in\_range.cpp}
distinct\_in\_ranges(a, queries): number of distinct values in a[l, r) per query. Sweeps r, keeps only last occurrences.
\begin{lstlisting}
vector<int> res = distinct_in_ranges(a, qs);  // qs[i] = {l, r}
\end{lstlisting}
\end{libentry}
\begin{libentry}{Dynamic (sparse) segment tree}{Ctrl+Alt+Y, D}{cpdynseg}{above solve()}{O(log range) per operation, O(q log range) memory}{library/ds/dynamic\_segtree.cpp}
DynSeg st(lo, hi): positions in [lo, hi) as large as 1e18, nodes made on demand; st.add(pos, v), st.sum(l, r) on [l, r).
\begin{lstlisting}
DynSeg st(0, 1e9);
st.add(x, 1); ll c = st.sum(l, r);
\end{lstlisting}
\end{libentry}
\begin{libentry}{Euler tour tree (dynamic forest: link, cut, connectivity, component sums)}{Ctrl+Alt+U, E}{cpett}{above solve()}{O(log n) expected per operation}{library/ds/euler\_tour\_tree.cpp}
ETT f(n); f.link(u, v), f.cut(u, v) (the edge must exist), f.connected(u, v), f.set(v, x), f.comp\_sum(v), f.comp\_size(v). Each tree is stored as its Euler tour in a treap.
\begin{lstlisting}
ETT f(n); f.link(0, 1); f.set(1, 5);
bool c = f.connected(0, 1); ll s = f.comp_sum(0); f.cut(0, 1);
\end{lstlisting}
\end{libentry}
\begin{libentry}{Fenwick tree / BIT}{Ctrl+Alt+T, I}{cpbit}{above solve()}{O(log n) per operation}{library/ds/fenwick.cpp}
BIT b(n); b.add(i, v), b.sum(l, r) on [l, r), b.prefix(i) = sum of [0, i), b.lower\_bound(k).
\begin{lstlisting}
BIT b(n);
b.add(i, v); b.sum(l, r); b.prefix(i); b.lower_bound(k);
\end{lstlisting}
\end{libentry}
\begin{libentry}{Fenwick tree 2D}{Ctrl+Alt+Y, Shift+I}{cpbit2d}{above solve()}{O(log n log m) per operation}{library/ds/fenwick\_2d.cpp}
BIT2D b(n, m); b.add(r, c, v); b.sum(r1, c1, r2, c2) = sum over rows [r1, r2) x cols [c1, c2).
\begin{lstlisting}
BIT2D b(n, m);
b.add(r, c, 1); ll s = b.sum(r1, c1, r2, c2);
\end{lstlisting}
\end{libentry}
\begin{libentry}{Fenwick tree: range add + range sum}{Ctrl+Alt+Y, I}{cpbitrange}{above solve()}{O(log n) per operation}{library/ds/fenwick\_range.cpp}
BITRange b(n); b.add(l, r, v) adds v on [l, r); b.sum(l, r). Two Fenwick trees over the difference array.
\begin{lstlisting}
BITRange b(n);
b.add(l, r, v); ll s = b.sum(l, r);
\end{lstlisting}
\end{libentry}
\begin{libentry}{Implicit treap (sequence with split / merge)}{Ctrl+Alt+Y, T}{cptreap}{above solve()}{O(log n) expected per operation}{library/ds/treap.cpp}
Treap tr; tr.insert(pos, v), tr.erase(pos), tr.reverse\_range(l, r), tr.range\_sum(l, r), tr.to\_vector(). Expected O(log n).
\begin{lstlisting}
Treap tr; for (auto x : a) tr.insert(tr.count(), x);
tr.reverse_range(l, r); ll s = tr.range_sum(l, r); vector<ll> now = tr.to_vector();
\end{lstlisting}
\end{libentry}
\begin{libentry}{Lazy segment tree (range add; range sum + min)}{Ctrl+Alt+T, Shift+S}{cplazy}{above solve()}{O(log n) per operation}{library/ds/lazy\_segment\_tree.cpp}
LazySeg st(n) or st(a); st.add(l, r, v), st.sum\_of(l, r), st.min\_of(l, r), all on [l, r).
\begin{lstlisting}
LazySeg st(a);
st.add(l, r, v); st.sum_of(l, r); st.min_of(l, r);  // all [l, r)
\end{lstlisting}
\end{libentry}
\begin{libentry}{Lazy segment tree (range assign + range add; sum, min)}{Ctrl+Alt+Y, A}{cplazyassign}{above solve()}{O(log n) per operation}{library/ds/lazy\_assign.cpp}
LazyAssign st(n or a): st.assign(l, r, v), st.add(l, r, v), st.sum\_of(l, r), st.min\_of(l, r), all on [l, r).
\begin{lstlisting}
LazyAssign st(a);
st.assign(l, r, 5); st.add(l, r, -2); st.sum_of(l, r); st.min_of(l, r);
\end{lstlisting}
\end{libentry}
\begin{libentry}{Mo\textquotesingle{}s algorithm (offline range queries)}{Ctrl+Alt+Y, O}{cpmo}{above solve()}{O((n + q) sqrt n) add/remove calls}{library/ds/mo\_algorithm.cpp}
mo(n, queries, add, remove, answer): queries [l, r); add(i) / remove(i) update the window, answer(qi) records query qi.
\begin{lstlisting}
vector<int> cnt(maxv), res(q); int distinct = 0;
mo(n, qs, [&](int i) { distinct += cnt[a[i]]++ == 0; }, [&](int i) { distinct -= --cnt[a[i]] == 0; }, [&](int qi) { res[qi] = distinct; });
\end{lstlisting}
\end{libentry}
\begin{libentry}{Mo\textquotesingle{}s algorithm on tree paths}{Ctrl+Alt+U, O}{cpmotree}{above solve()}{O((n + q) sqrt(n) * cost of add / remove + q log n)}{library/ds/mo\_tree.cpp}
mo\_on\_tree(adj, root, qs, add, remove, answer): offline queries on the vertex set of the path u-v. add(v)/remove(v) change the current set, answer(i) records query i. Uses the Euler tour trick (a vertex seen twice cancels).
\begin{lstlisting}
mo_on_tree(adj, 0, qs, [&](int v) { if (cnt[c[v]]++ == 0) distinct++; }, [&](int v) { if (--cnt[c[v]] == 0) distinct--; }, [&](int i) { ans[i] = distinct; });
\end{lstlisting}
\end{libentry}
\begin{libentry}{Mo\textquotesingle{}s algorithm with updates (time as a third dimension)}{Ctrl+Alt+U, Shift+O}{cpmoupdates}{above solve()}{O(n\textasciicircum{}(5/3)) for n queries and updates}{library/ds/mo\_updates.cpp}
mo\_with\_updates(n, qs, add, remove, apply, answer): qs[i] = \{l, r, t\} asks about [l, r) after the first t updates. apply(u, L, R) toggles update u (swap in/out) given the window [L, R).
\begin{lstlisting}
mo_with_updates(n, qs, add, remove, [&](int u, int L, int R) { auto& [p, v] = upd[u]; if (L <= p && p < R) remove(p); swap(a[p], v); if (L <= p && p < R) add(p); }, [&](int i) { ans[i] = cur; });
\end{lstlisting}
\end{libentry}
\begin{libentry}{Ordered set (PBDS order statistics)}{Ctrl+Alt+T, Shift+O}{cposet}{above solve()}{O(log n) per operation}{library/ds/ordered\_set.cpp}
ordered\_set\textless{}int\textgreater{} s; s.find\_by\_order(k) = k-th smallest (0-based), s.order\_of\_key(x) = \# elements \textless{} x.
\begin{lstlisting}
ordered_set<int> s;
s.insert(x); *s.find_by_order(k); s.order_of_key(x);
\end{lstlisting}
\end{libentry}
\begin{libentry}{Persistent implicit treap (rope with versions)}{Ctrl+Alt+U, R}{cpptreap}{above solve()}{O(log n) expected per operation, O(log n) new nodes per update}{library/ds/persistent\_treap.cpp}
PTreap t; int v = t.build(a); every operation returns a NEW version and keeps the old: insert(v, pos, x), erase(v, pos), concat(v1, v2), slice(v, l, r). Queries: get(v, pos), sum(v, l, r), size(v).
\begin{lstlisting}
PTreap t; int v0 = t.build(a); int v1 = t.insert(v0, 2, 7);
int copy_paste = t.concat(v1, t.slice(v1, l, r)); ll s = t.sum(v0, l, r);
t.t.reserve(n + 50 * q);  // ~46 new 32-byte nodes per update; regrowing briefly holds two copies
\end{lstlisting}
\end{libentry}
\begin{libentry}{Persistent segment tree (k-th smallest in range)}{Ctrl+Alt+Y, P}{cppersist}{above solve()}{O(log n) per operation and O(log n) new nodes per update}{library/ds/persistent\_segtree.cpp}
PersistentSeg ps(n): root = ps.build(a) in O(n) nodes; root2 = ps.update(root1, pos, v) keeps every version; ps.query(root, l, r); ps.kth(rootL, rootR, k).
\begin{lstlisting}
vector<int> root{0}; for (int i = 0; i < n; i++) root.push_back(ps.update(root[i], rank[i], 1));  // ranks from cpcompress
int idx = ps.kth(root[l], root[r], k);  // k-th smallest (0-based) of a[l, r) is vals[idx]
int r0 = ps.build(a);  // a version holding the values a, instead of n updates (n log n nodes)
\end{lstlisting}
\end{libentry}
\begin{libentry}{Segment tree (point update, range query)}{Ctrl+Alt+T, S}{cpseg}{above solve()}{O(log n) per operation}{library/ds/segment\_tree.cpp}
SegTree st(n) or st(a); st.set(i, x), st.get(i), st.query(l, r) on [l, r). Edit T / ID / op.
\begin{lstlisting}
SegTree st(a);  // or SegTree st(n)
st.set(i, x); st.get(i); st.query(l, r);  // [l, r)
\end{lstlisting}
\end{libentry}
\begin{libentry}{Segment tree beats (range chmin, sum, max)}{Ctrl+Alt+Y, S}{cpbeats}{above solve()}{O((n + q) log\textasciicircum{}2 n) amortized}{library/ds/segtree\_beats.cpp}
Beats st(a): st.chmin(l, r, x) sets a[i] = min(a[i], x) on [l, r); st.sum(l, r), st.max(l, r). Amortized O(log\textasciicircum{}2 n).
\begin{lstlisting}
Beats st(a);
st.chmin(l, r, x); ll s = st.sum(l, r); ll m = st.max(l, r);
\end{lstlisting}
\end{libentry}
\begin{libentry}{Segment tree with custom nodes + binary search}{Ctrl+Alt+Y, N}{cpsegnode}{above solve()}{O(log n) per operation}{library/ds/segtree\_node.cpp}
SegNode st(n or leaves): any associative Node (example: max subarray sum); st.max\_right(l, pred) walks down in O(log n).
\begin{lstlisting}
SegNode st(n); st.set(i, Node::leaf(x)); Node r = st.query(l, r); r.best;
int r = st.max_right(l, [&](const Node& x) { return x.sum < k; });  // first r with sum of [l, r] >= k is r
\end{lstlisting}
\end{libentry}
\begin{libentry}{Segment tree: range add, min + count of min}{Ctrl+Alt+Y, Shift+C}{cpmincount}{above solve()}{O(log n) per operation}{library/ds/min\_count\_segtree.cpp}
MinCount st(weights): add(l, r, v); query(l, r) = \{min, total weight at the min\}; zero\_weight() = weight of cells equal to 0 (values must stay \textgreater{}= 0, as coverage counts do).
\begin{lstlisting}
MinCount st(len);  // len[i] = width of elementary interval i (all 1 for plain counting)
st.add(l, r, +1); st.add(l, r, -1); ll covered = total_width - st.zero_weight();  // union of rectangles sweep
\end{lstlisting}
\end{libentry}
\begin{libentry}{Sparse table (O(1) range min)}{Ctrl+Alt+T, R}{cpsparse}{above solve()}{O(n log n) build, O(1) per query}{library/ds/sparse\_table.cpp}
SparseTable st(a); st.query(l, r) = min of [l, r), l \textless{} r. Swap op for max/gcd (idempotent only).
\begin{lstlisting}
SparseTable st(a);
st.query(l, r);  // min of [l, r)
\end{lstlisting}
\end{libentry}
\begin{libentry}{Wavelet tree (k-th smallest / count \textless{}= x in range)}{Ctrl+Alt+Y, W}{cpwavelet}{above solve()}{O(n log n) build, O(log n) per query}{library/ds/wavelet\_tree.cpp}
Wavelet wt(a): wt.kth(l, r, k) = k-th smallest (0-based) of a[l, r); wt.count\_leq(l, r, x). Any ll values (compressed inside); a bit-packed wavelet matrix, so all levels stay in cache.
\begin{lstlisting}
Wavelet wt(a);
ll med = wt.kth(l, r, (r - l) / 2); int c = wt.count_leq(l, r, x);
\end{lstlisting}
\end{libentry}
\section{Dynamic programming}
\begin{libentry}{(min, +) convolution with convex sequences}{Ctrl+Alt+U, Shift+C}{cpminplus}{above solve()}{O(n + m) / O((n + m) log(n + m))}{library/dp/minplus\_convex.cpp}
c[k] = min over i + j = k of a[i] + b[j]. minplus\_convex\_convex(a, b): both convex, O(n + m) by merging slopes; minplus\_convex\_any(a, b): a convex, b arbitrary, O((n + m) log).
\begin{lstlisting}
vector<ll> c = minplus_convex_convex(a, b);
vector<ll> d = minplus_convex_any(convex_a, arbitrary_b);
\end{lstlisting}
\end{libentry}
\begin{libentry}{Bitmask DP (travelling salesman)}{Ctrl+Alt+T, ]}{cptsp}{above solve()}{O(2\textasciicircum{}n n\textasciicircum{}2)}{library/dp/bitmask\_tsp.cpp}
tsp(dist): shortest cycle through all n \textless{}= \textasciitilde{}18 vertices from 0. dp[mask][v] = best path covering mask, ending at v.
\begin{lstlisting}
ll best = tsp(dist);  // dist[i][j] = cost i -> j
// Hamiltonian path count: same loops with dp[mask | 1 << v][v] += dp[mask][u] over edges
\end{lstlisting}
\end{libentry}
\begin{libentry}{Broken profile DP (domino tilings)}{Ctrl+Alt+Y, Shift+P}{cpprofile}{above solve()}{O(n m 2\textasciicircum{}m)}{library/dp/broken\_profile.cpp}
domino\_tilings(n, m): ways to tile an n x m board with 1 x 2 dominoes mod MOD, cell by cell with an m-bit profile (m \textless{}= \textasciitilde{}12).
\begin{lstlisting}
ll ways = domino_tilings(n, m);  // put the smaller side as m
\end{lstlisting}
\end{libentry}
\begin{libentry}{Digit DP}{Ctrl+Alt+T, =}{cpdigit}{above solve()}{O(digits * states * 10)}{library/dp/digit\_dp.cpp}
count\_upto(N) = how many x in [0, N] have no two equal adjacent digits. Change the state / check for your property.
\begin{lstlisting}
ll ans = count_upto(B) - count_upto(A - 1);  // count in [A, B]
\end{lstlisting}
\end{libentry}
\begin{libentry}{Divide \& conquer DP optimization}{Ctrl+Alt+Y, Q}{cpdnc}{above solve()}{O(K n log n) cost calls}{library/dp/dnc\_optimization.cpp}
dnc\_opt(K, n, cost): min total cost splitting a[0, n) into K groups; cost(j, i) = cost of group a[j, i). Needs monotone optimal split points.
\begin{lstlisting}
ll best = dnc_opt(K, n, [&](int j, int i) { ll s = pre[i] - pre[j]; return s * s; });
\end{lstlisting}
\end{libentry}
\begin{libentry}{Grid paths DP}{Ctrl+Alt+T, Shift+Z}{cpgridpaths}{inside solve()}{O(nm)}{library/dp/grid\_paths.cpp}
needs n, m, grid. ways[r][c] = right/down paths from (0, 0) to (r, c) avoiding \textquotesingle{}\#\textquotesingle{}, mod MOD.
\begin{lstlisting}
// answer: ways[n - 1][m - 1]; for a max-sum path use max(...) + value instead of the sum
\end{lstlisting}
\end{libentry}
\begin{libentry}{Knapsack (0/1, unbounded, coin-change ways)}{Ctrl+Alt+T, Shift+Q}{cpknap}{above solve()}{O(n W)}{library/dp/knapsack.cpp}
knapsack\_01(items, W) / knapsack\_unbounded(items, W): max value, weight \textless{}= W; coin\_ways(coins, X): combinations mod MOD.
\begin{lstlisting}
vector<pair<int, ll>> items;  // {weight, value}
ll best = knapsack_01(items, W); ll ways = coin_ways(coins, X);
\end{lstlisting}
\end{libentry}
\begin{libentry}{Knuth optimization (interval DP)}{Ctrl+Alt+Y, Shift+Q}{cpknuth}{above solve()}{O(n\textasciicircum{}2)}{library/dp/knuth\_optimization.cpp}
knuth\_opt(n, cost): dp[l][r] = min over l \textless{} m \textless{} r of dp[l][m] + dp[m][r] + cost(l, r), dp[i][i+1] = 0. Needs monotone opt.
\begin{lstlisting}
ll best = knuth_opt(n, [&](int l, int r) { return pre[r] - pre[l]; });  // optimal merging of a[0, n)
\end{lstlisting}
\end{libentry}
\begin{libentry}{LCS + edit distance}{Ctrl+Alt+T, [}{cplcs}{above solve()}{O(nm)}{library/dp/lcs\_edit.cpp}
lcs(a, b) -\textgreater{} \{length, one LCS\}; edit\_distance(a, b) = min insert / delete / replace operations.
\begin{lstlisting}
auto [len, s] = lcs(a, b);
int d = edit_distance(a, b);
\end{lstlisting}
\end{libentry}
\begin{libentry}{Li Chao tree (min of lines)}{Ctrl+Alt+Y, Shift+L}{cplichao}{above solve()}{O(log range) per add / query}{library/dp/li\_chao.cpp}
LiChao lc(lo, hi); lc.add(k, m); lc.query(x) = min of k*x + m for integer x in [lo, hi). Nodes created on demand.
\begin{lstlisting}
LiChao lc(-1e9, 1e9 + 1);
lc.add(k, m); ll best = lc.query(x);
\end{lstlisting}
\end{libentry}
\begin{libentry}{LineContainer (convex hull trick, any order)}{Ctrl+Alt+Y, L}{cplinecontainer}{above solve()}{O(log n) per add / query}{library/dp/line\_container.cpp}
LineContainer lc; lc.add(k, m); lc.query(x) = max of k*x + m. For a minimum add(-k, -m) and negate the query.
\begin{lstlisting}
LineContainer lc;
lc.add(slope, intercept); ll best = lc.query(x);  // dp[i] = max_j (dp[j] + b[j] * a[i]) in O(log n)
\end{lstlisting}
\end{libentry}
\begin{libentry}{Longest increasing subsequence}{Ctrl+Alt+T, Shift+L}{cplis}{above solve()}{O(n log n)}{library/dp/lis.cpp}
vector\textless{}int\textgreater{} idx = lis(a): indices of one strictly increasing LIS, O(n log n). Length = idx.size().
\begin{lstlisting}
vector<int> idx = lis(a);  // idx.size() = LIS length
\end{lstlisting}
\end{libentry}
\begin{libentry}{Monotone convex hull trick (sorted slopes)}{Ctrl+Alt+Y, Z}{cpmonocht}{above solve()}{O(1) amortized add, O(log n) query, O(1) amortized query\_inc}{library/dp/monotone\_cht.cpp}
MonoCHT h; h.add(k, m) with non-increasing k; h.query(x) = min of k*x + m (binary search); h.query\_inc(x) for non-decreasing x (pointer walk). For max: add(-k, -m), negate.
\begin{lstlisting}
MonoCHT h; h.add(b[j], dp[j]);
dp[i] = h.query_inc(a[i]) + c[i];
\end{lstlisting}
\end{libentry}
\begin{libentry}{Monotone minima (row minima when the argmin moves right)}{Ctrl+Alt+U, M}{cpmonominima}{above solve()}{O((n + m) log n) evaluations of f}{library/dp/monotone\_minima.cpp}
monotone\_minima(n, m, f): for each row i, the leftmost column j minimizing f(i, j), given that this argmin is non-decreasing in i (e.g. a Monge matrix).
\begin{lstlisting}
vector<int> opt = monotone_minima(n, m, [&](int i, int j) { return cost(i, j); });
\end{lstlisting}
\end{libentry}
\begin{libentry}{Range (interval) DP}{Ctrl+Alt+T, -}{cprangedp}{inside solve()}{O(n\textasciicircum{}3)}{library/dp/range\_dp.cpp}
needs n, a. dp[l][r] over intervals by increasing length; example: merge adjacent piles, cost = merged sum.
\begin{lstlisting}
// answer: dp[0][n - 1]; change the inner transition for your problem
\end{lstlisting}
\end{libentry}
\begin{libentry}{Slope trick (make an array non-decreasing)}{Ctrl+Alt+Y, Shift+S}{cpslope}{above solve()}{O(n log n)}{library/dp/slope\_trick.cpp}
make\_non\_decreasing(a): min total \textbar{}a[i] - b[i]\textbar{} over non-decreasing b. Max-heap of slope breakpoints. Strictly increasing: use a[i] - i.
\begin{lstlisting}
ll cost = make_non_decreasing(a);
\end{lstlisting}
\end{libentry}
\begin{libentry}{Subset convolution}{Ctrl+Alt+U, Shift+S}{cpsubsetconv}{above solve()}{O(n\textasciicircum{}2 2\textasciicircum{}n)}{library/dp/subset\_convolution.cpp}
subset\_convolution(a, b): c[S] = sum over T subset of S of a[T] * b[S \textbackslash{} T] (disjoint unions only), mod MOD. Sizes 2\textasciicircum{}n; n = 20 takes about 1 s and 176 MB.
\begin{lstlisting}
vector<ll> c = subset_convolution(a, b);   // a.size() == b.size() == 1 << n
\end{lstlisting}
\end{libentry}
\begin{libentry}{Subset sums with bitset}{Ctrl+Alt+T, Shift+1}{cpbitsetsum}{above solve()}{O(n MAXS / 64)}{library/dp/subset\_sum\_bitset.cpp}
subset\_sums\textless{}MAXS\textgreater{}(a): bitset with bit s set iff some subset of a sums to s (s \textless{} MAXS). 64x faster than a bool DP.
\begin{lstlisting}
auto can = subset_sums<100001>(a);
if (can[target]) { ... }
\end{lstlisting}
\end{libentry}
\begin{libentry}{Sum over subsets (SOS DP)}{Ctrl+Alt+Y, U}{cpsos}{above solve()}{O(B 2\textasciicircum{}B)}{library/dp/sos\_dp.cpp}
sos\_subsets(f, B): f[mask] becomes the sum of f[sub] over all submasks; sos\_supersets: over supermasks. In place, size 2\textasciicircum{}B.
\begin{lstlisting}
vector<ll> f(1 << B); /* f[mask] = count of numbers equal to mask */ sos_subsets(f, B);
// f[mask] now counts numbers x with (x & mask) == x
\end{lstlisting}
\end{libentry}
\section{Math}
\begin{libentry}{Berlekamp-Massey (find a linear recurrence)}{Ctrl+Alt+U, Shift+B}{cpbm}{above solve()}{O(n\textasciicircum{}2)}{library/math/berlekamp\_massey.cpp}
berlekamp\_massey(s): shortest c with s[i] = c[0]*s[i-1] + c[1]*s[i-2] + ... mod MOD (prime). Feed 2L terms to find a recurrence of order L.
\begin{lstlisting}
vector<ll> c = berlekamp_massey({1, 1, 2, 3, 5, 8, 13, 21});   // {1, 1}
ll x = linear_rec_nth(first_terms, c, n);   // with the Kitamasa template
\end{lstlisting}
\end{libentry}
\begin{libentry}{Chinese remainder theorem (any moduli)}{Ctrl+Alt+Y, 1}{cpcrt}{above solve()}{O(log m)}{library/math/crt.cpp}
crt(r1, m1, r2, m2) -\textgreater{} \{r, lcm\}: x = r (mod lcm) solves both congruences; \{-1, -1\} if none. Moduli need not be coprime.
\begin{lstlisting}
auto [r, m] = crt(2, 3, 3, 5);  // x = 8 (mod 15)
// many congruences: fold left, feeding each result back in as (r1, m1); lcm must stay < ~4e18
\end{lstlisting}
\end{libentry}
\begin{libentry}{Determinant mod p, matrix-tree theorem}{Ctrl+Alt+U, Shift+D}{cpdet}{above solve()}{O(n\textasciicircum{}3)}{library/math/determinant.cpp}
det\_mod(A): determinant mod MOD (prime) by Gaussian elimination; count\_spanning\_trees(n, edges): Kirchhoff\textquotesingle{}s matrix-tree theorem (multi-edges allowed). LGV lemma: det of path-count matrix.
\begin{lstlisting}
ll d = det_mod(A);
ll t = count_spanning_trees(n, edges);   // edges: 0-indexed pairs
\end{lstlisting}
\end{libentry}
\begin{libentry}{Discrete logarithm (baby-step giant-step)}{Ctrl+Alt+U, L}{cpdlog}{above solve()}{O(sqrt(m)) time and memory}{library/math/discrete\_log.cpp}
discrete\_log(a, b, m): smallest x \textgreater{}= 0 with a\textasciicircum{}x = b (mod m), or -1. Any m (not only primes), m \textless{}= 3e9.
\begin{lstlisting}
ll x = discrete_log(2, 5, 13);   // 2^9 = 512 = 5 (mod 13) -> 9
\end{lstlisting}
\end{libentry}
\begin{libentry}{Divisor sieve (count + sum of divisors)}{Ctrl+Alt+T, Shift+2}{cpdivsieve}{above solve()}{O(N log N)}{library/math/divisor\_sieve.cpp}
divisor\_sieve(N) -\textgreater{} \{cnt, sum\}: number and sum of divisors of every x \textless{}= N by looping over multiples.
\begin{lstlisting}
auto [cnt, sum] = divisor_sieve(1000000);
\end{lstlisting}
\end{libentry}
\begin{libentry}{Du\textquotesingle{}s sieve: prefix sums of phi and mu up to 1e10}{Ctrl+Alt+U, D}{cpdusieve}{above solve()}{O(n\textasciicircum{}(2/3)) for the first query, memoised after}{library/math/du\_sieve.cpp}
DuSieve ds(n\_max); ds.mu\_sum(n) = mu(1)+...+mu(n) (exact), ds.phi\_sum(n) = phi(1)+...+phi(n) mod MOD.
\begin{lstlisting}
DuSieve ds(10'000'000'000LL);
ll M = ds.mu_sum(n), Phi = ds.phi_sum(n);
\end{lstlisting}
\end{libentry}
\begin{libentry}{Euler\textquotesingle{}s totient}{Ctrl+Alt+Y, Shift+E}{cpphi}{above solve()}{O(sqrt n) single, O(N log log N) sieve}{library/math/totient.cpp}
phi(n) by trial division; phi\_sieve(N) = phi of every x \textless{}= N. phi(n) counts 1 \textless{}= k \textless{}= n with gcd(k, n) = 1.
\begin{lstlisting}
ll p = phi(n);
auto ph = phi_sieve(1000000);  // a^phi(m) = 1 (mod m) when gcd(a, m) = 1
\end{lstlisting}
\end{libentry}
\begin{libentry}{Extended GCD + inverse for any modulus}{Ctrl+Alt+T, X}{cpextgcd}{above solve()}{O(log min(a, b))}{library/math/ext\_gcd.cpp}
g = ext\_gcd(a, b, x, y) with a*x + b*y = g; inv\_mod(a, m) for gcd(a, m) = 1, else -1.
\begin{lstlisting}
ll x, y, g = ext_gcd(a, b, x, y);  // a*x + b*y = g
inv_mod(a, m);  // -1 if gcd(a, m) != 1
\end{lstlisting}
\end{libentry}
\begin{libentry}{Factorize + divisors (trial division)}{Ctrl+Alt+T, Shift+P}{cpfactor}{above solve()}{O(sqrt x)}{library/math/factorize.cpp}
factorize(x) -\textgreater{} \{prime, exponent\}, divisors(x) sorted. O(sqrt x), fine up to \textasciitilde{}1e12.
\begin{lstlisting}
for (auto [p, e] : factorize(x)) { ... }
vector<ll> d = divisors(x);
\end{lstlisting}
\end{libentry}
\begin{libentry}{FFT (exact integer convolution)}{Ctrl+Alt+Y, F}{cpfft}{above solve()}{O((n + m) log(n + m))}{library/math/fft.cpp}
multiply(a, b): convolution c[k] = sum a[i] b[k - i] with complex FFT; exact while results stay below \textasciitilde{}1e15.
\begin{lstlisting}
vector<ll> c = multiply(a, b);  // polynomial product, counting pairs with a given sum, ...
\end{lstlisting}
\end{libentry}
\begin{libentry}{Floor-division blocks}{Ctrl+Alt+Y, Shift+F}{cpfloorblocks}{above solve()}{O(sqrt n)}{library/math/floor\_blocks.cpp}
for\_each\_floor\_block(n, f): f(l, r, q) for every maximal [l, r] where n / i == q; only O(sqrt n) blocks.
\begin{lstlisting}
ll s = 0; for_each_floor_block(n, [&](ll l, ll r, ll q) { s += q * (r - l + 1); });  // sum of n / i
// sum of divisors of 1..n: q * (l + r) * (r - l + 1) / 2 per block
\end{lstlisting}
\end{libentry}
\begin{libentry}{Gaussian elimination mod p (linear systems, rank)}{Ctrl+Alt+U, N}{cpgauss}{above solve()}{O(n m min(n, m))}{library/math/gauss\_mod.cpp}
gauss\_mod(a, b, x): solves a x = b (mod MOD, MOD prime) for an n x m matrix a. Returns the rank, or -1 if there is no solution; x receives one solution (free variables set to 0).
\begin{lstlisting}
vector<vector<ll>> a(n, vector<ll>(m)); vector<ll> b(n), x;
int r = gauss_mod(a, b, x);   // -1: inconsistent; else m - r free variables, MOD^(m - r) solutions
\end{lstlisting}
\end{libentry}
\begin{libentry}{Grundy numbers (Sprague-Grundy)}{Ctrl+Alt+Y, G}{cpgrundy}{above solve()}{O(N * moves)}{library/math/grundy.cpp}
grundy\_numbers(N, moves): g[s] = mex of g over moves(s) (moves must go to smaller states). XOR of games != 0 -\textgreater{} first player wins.
\begin{lstlisting}
auto g = grundy_numbers(N, [&](int s) { vector<int> to; for (int k : {1, 3, 4}) if (s >= k) to.push_back(s - k); return to; });
int x = 0; for (int pile : piles) x ^= g[pile]; bool first_wins = x != 0;
\end{lstlisting}
\end{libentry}
\begin{libentry}{Inclusion-exclusion (divisible by any)}{Ctrl+Alt+Y, 3}{cpincexc}{above solve()}{O(2\textasciicircum{}k k)}{library/math/inclusion\_exclusion.cpp}
count\_divisible\_by\_any(N, ds): how many x in [1, N] are divisible by at least one of ds (k \textless{}= \textasciitilde{}20). Odd subsets add, even subtract.
\begin{lstlisting}
ll c = count_divisible_by_any(N, {2, 3, 5});
// coprime to m: N - count_divisible_by_any(N, distinct primes of m)
\end{lstlisting}
\end{libentry}
\begin{libentry}{Lagrange interpolation}{Ctrl+Alt+U, I}{cplagrange}{above solve()}{O(k) consecutive points (plus O(log MOD)), O(k\textasciicircum{}2) general points}{library/math/lagrange.cpp}
lagrange\_consecutive(y, x): the degree \textless{} k polynomial with P(i) = y[i] for i = 0..k-1, evaluated at any x; lagrange\_general(xs, ys, x) for arbitrary distinct points. Mod MOD (prime).
\begin{lstlisting}
ll s = lagrange_consecutive({0, 1, 5, 14}, 1'000'000'000);   // 0^2 + 1^2 + ... + 10^9^2 mod MOD
ll v = lagrange_general(xs, ys, x);
\end{lstlisting}
\end{libentry}
\begin{libentry}{Matrix exponentiation}{Ctrl+Alt+T, Shift+X}{cpmatrix}{above solve()}{O(k\textasciicircum{}3 log e)}{library/math/matrix.cpp}
Mat R = mat\_pow(A, e), mat\_mul(A, B), entries mod MOD. Linear recurrences in O(k\textasciicircum{}3 log e).
\begin{lstlisting}
Mat M = {{1, 1}, {1, 0}};
ll fib_n = mat_pow(M, n)[0][1];
\end{lstlisting}
\end{libentry}
\begin{libentry}{Miller-Rabin + Pollard rho (numbers up to 1e18)}{Ctrl+Alt+Y, X}{cppollard}{above solve()}{\textasciitilde{}O(n\textasciicircum{}(1/4) log n) per factorization}{library/math/pollard\_rho.cpp}
BigFactor::is\_prime(n) (deterministic below 2\textasciicircum{}64), BigFactor::factor(n) -\textgreater{} prime factors with repeats, unsorted.
\begin{lstlisting}
bool p = BigFactor::is_prime(n);
auto f = BigFactor::factor(n); sort(f.begin(), f.end());
\end{lstlisting}
\end{libentry}
\begin{libentry}{Mint (modular integer)}{Ctrl+Alt+T, Shift+M}{cpmint}{above solve()}{O(1) per op, O(log MOD) for / and inv}{library/math/mint.cpp}
mint x = 5; x += 3; x *= x.inv(); x.pow(e); cout \textless{}\textless{} x. Uses the template\textquotesingle{}s MOD (must be prime for /).
\begin{lstlisting}
mint x = 5;
x = x * 3 - 1; x /= 2; x.pow(10); cout << x;
\end{lstlisting}
\end{libentry}
\begin{libentry}{Mobius function sieve}{Ctrl+Alt+Y, Shift+M}{cpmobius}{above solve()}{O(N)}{library/math/mobius.cpp}
mobius\_sieve(N) = mu[0..N] by a linear sieve. Key identity: sum over d \textbar{} g of mu(d) = [g == 1] (count coprime pairs, etc.).
\begin{lstlisting}
auto mu = mobius_sieve(N);
// pairs (a, b) <= N with gcd 1: sum over d of mu[d] * (N / d) * (N / d)
\end{lstlisting}
\end{libentry}
\begin{libentry}{Modular power + inverse}{Ctrl+Alt+T, M}{cpmodpow}{above solve()}{O(log e)}{library/math/modpow.cpp}
modpow(a, e, m = MOD), modinv(a, m = MOD) for prime m. Uses the template\textquotesingle{}s MOD.
\begin{lstlisting}
modpow(2, n);  // 2^n mod MOD
modinv(x);    // x^-1 mod MOD
\end{lstlisting}
\end{libentry}
\begin{libentry}{N-th term of a linear recurrence (Kitamasa)}{Ctrl+Alt+U, K}{cpkitamasa}{above solve()}{O(L\textasciicircum{}2 log n)}{library/math/linear\_recurrence.cpp}
linear\_rec\_nth(a, c, n): a[n] for a[i] = c[0]*a[i-1] + ... + c[L-1]*a[i-L], given a[0..L-1]; mod MOD. Pairs with Berlekamp-Massey.
\begin{lstlisting}
ll f = linear_rec_nth({0, 1}, {1, 1}, 1'000'000'000'000'000'000LL);   // Fibonacci
\end{lstlisting}
\end{libentry}
\begin{libentry}{nCr / factorials mod p}{Ctrl+Alt+T, N}{cpncr}{above solve()}{O(N) init, O(1) per query}{library/math/ncr.cpp}
init\_comb(N) once, then nCr, nPr, catalan, stars\_and\_bars, fact[i], inv\_fact[i]. Uses the template\textquotesingle{}s MOD.
\begin{lstlisting}
init_comb(200000);
nCr(n, r); nPr(n, r); catalan(n); fact[n];
\end{lstlisting}
\end{libentry}
\begin{libentry}{NTT (convolution mod 998244353)}{Ctrl+Alt+Y, Shift+T}{cpntt}{above solve()}{O((n + m) log(n + m))}{library/math/ntt.cpp}
ntt\_multiply(a, b): convolution with every coefficient mod 998244353, exact for any size up to 2\textasciicircum{}23.
\begin{lstlisting}
vector<ll> c = ntt_multiply(a, b);
\end{lstlisting}
\end{libentry}
\begin{libentry}{Pascal triangle (nCr for any modulus)}{Ctrl+Alt+Y, Shift+N}{cppascal}{above solve()}{O(N\textasciicircum{}2)}{library/math/pascal.cpp}
pascal(N, m)[n][r] = C(n, r) mod m for n \textless{}= N; works when m is not prime or is 1e18-ish (then keep N small).
\begin{lstlisting}
auto C = pascal(60, 1e18);  // exact binomials up to C(60, 30) ~ 1.2e17
\end{lstlisting}
\end{libentry}
\begin{libentry}{Polynomial toolkit mod 998244353 (inverse, divmod, log, exp, sqrt, pow, multipoint)}{Ctrl+Alt+U, Shift+P}{cppoly}{above solve()}{O(n log n) for mul / inv / divmod / log / exp / sqrt / pow; O(n log\textasciicircum{}2 n) for multipoint and interpolation}{library/math/poly.cpp}
Poly = vector\textless{}ll\textgreater{}. poly\_mul, poly\_inv(a, n), poly\_divmod(a, b), poly\_log / poly\_exp / poly\_sqrt (a, n), poly\_pow(a, k, n), poly\_eval\_many(a, xs), poly\_interpolate(xs, ys). Self-contained NTT.
\begin{lstlisting}
Poly b = poly_inv(a, n);   // a * b = 1 mod x^n, needs a[0] != 0
auto [q, r] = poly_divmod(a, b);  Poly e = poly_exp(a, n);  vector<ll> v = poly_eval_many(a, xs);
\end{lstlisting}
\end{libentry}
\begin{libentry}{Prime counting and prime sums up to 1e11 (Lucy)}{Ctrl+Alt+U, P}{cpprimecount}{above solve()}{O(n\textasciicircum{}(3/4)) time, O(sqrt n) memory}{library/math/prime\_count.cpp}
prime\_count(n) = number of primes \textless{}= n (exact); prime\_sum(n) = sum of primes \textless{}= n mod MOD.
\begin{lstlisting}
ll pi = prime_count(100'000'000'000LL);   // 4118054813
ll s = prime_sum(n);
\end{lstlisting}
\end{libentry}
\begin{libentry}{Primitive root}{Ctrl+Alt+U, G}{cpprimroot}{above solve()}{O(sqrt(m) + answer * (number of prime factors of phi) * log m)}{library/math/primitive\_root.cpp}
primitive\_root(m): smallest generator g of (Z/mZ)* (exists for m = 1, 2, 4, p\textasciicircum{}k, 2p\textasciicircum{}k), else -1. m \textless{}= 3e9.
\begin{lstlisting}
ll g = primitive_root(998244353);   // 3
\end{lstlisting}
\end{libentry}
\begin{libentry}{Sieve (primes + smallest prime factor)}{Ctrl+Alt+T, P}{cpsieve}{above solve()}{O(N) init, O(log x) per factorization}{library/math/sieve.cpp}
init\_sieve(N); primes, spf[x], is\_prime(x), factor\_small(x) -\textgreater{} \{p, e\} in O(log x) for x \textless{}= N.
\begin{lstlisting}
init_sieve(10000000);
is_prime(x); primes; factor_small(x);  // x <= N
\end{lstlisting}
\end{libentry}
\begin{libentry}{Square root modulo a prime (Tonelli-Shanks)}{Ctrl+Alt+U, Q}{cpsqrtmod}{above solve()}{O(log\textasciicircum{}2 p)}{library/math/sqrt\_mod.cpp}
sqrt\_mod(a, p): some x with x*x = a (mod p) for prime p (the other root is p - x), or -1 if a is not a square.
\begin{lstlisting}
ll x = sqrt_mod(10, 13);   // 6 or 7: 6*6 = 36 = 10 (mod 13)
\end{lstlisting}
\end{libentry}
\begin{libentry}{Stern-Brocot: simplest fraction in an interval, best approximation}{Ctrl+Alt+U, B}{cpsternbrocot}{above solve()}{O(log max) (a Euclid-like walk down the Stern-Brocot tree)}{library/math/stern\_brocot.cpp}
simplest\_between(a, b, c, d): p/q with the smallest q (then p) and a/b \textless{} p/q \textless{} c/d (0 \textless{}= a/b \textless{} c/d); best\_approx(x, y, N): p/q closest to x/y with q \textless{}= N.
\begin{lstlisting}
auto [p, q] = simplest_between(1, 3, 1, 2);   // 2/5
auto [p2, q2] = best_approx(314159, 100000, 100);   // 311/99
\end{lstlisting}
\end{libentry}
\begin{libentry}{Stirling numbers (both kinds) and Bell numbers}{Ctrl+Alt+U, S}{cpstirling}{above solve()}{O(N\textasciicircum{}2) time and memory}{library/math/stirling.cpp}
init\_stirling(N): S2[n][k] (partitions of n items into k nonempty sets), S1[n][k] (permutations of n with k cycles), bell[n]; all mod MOD.
\begin{lstlisting}
init_stirling(1000);
S2[n][k]; S1[n][k]; bell[n];
\end{lstlisting}
\end{libentry}
\begin{libentry}{Walsh-Hadamard transform: XOR / AND / OR convolution}{Ctrl+Alt+U, W}{cpfwht}{above solve()}{O(n log n) for n = padded size}{library/math/fwht.cpp}
xor\_convolution(a, b), and\_convolution(a, b), or\_convolution(a, b): c[k] = sum of a[i]*b[j] over i\textasciicircum{}j = k (i\&j, i\textbar{}j) mod MOD. Sizes padded to a power of two.
\begin{lstlisting}
vector<ll> c = xor_convolution(a, b);
\end{lstlisting}
\end{libentry}
\begin{libentry}{XOR basis (linear basis over GF(2))}{Ctrl+Alt+Y, Shift+X}{cpxorbasis}{above solve()}{O(B) per operation, B = 62 bits}{library/math/xor\_basis.cpp}
XorBasis xb; xb.insert(x) (false if x was already representable); xb.can\_make(x); xb.max\_xor(); 2\textasciicircum{}xb.dim distinct subset XORs.
\begin{lstlisting}
XorBasis xb; for (auto x : a) xb.insert(x);
ll best = xb.max_xor(); bool ok = xb.can_make(target);
\end{lstlisting}
\end{libentry}
\section{Strings}
\begin{libentry}{Aho-Corasick (many patterns at once)}{Ctrl+Alt+Y, Shift+A}{cpaho}{above solve()}{O(total pattern length * 26 + \textbar{}text\textbar{})}{library/string/aho\_corasick.cpp}
AhoCorasick ac; id = ac.add(p) per pattern; ac.build(); ac.occurrences(text)[id] = matches of that pattern. a-z only.
\begin{lstlisting}
AhoCorasick ac; vector<int> id; for (auto& p : pats) id.push_back(ac.add(p)); ac.build();
auto occ = ac.occurrences(text); ll hits = occ[id[i]];  // or walk: s = ac.go(s, ch), ac.out[s] = matches ending here
\end{lstlisting}
\end{libentry}
\begin{libentry}{Binary trie (maximum XOR)}{Ctrl+Alt+T, \textasciigrave{}}{cpxortrie}{above solve()}{O(B) per operation, B = 30 (set B = 61 for 64-bit values)}{library/string/xor\_trie.cpp}
XorTrie t; t.insert(x), t.erase(x) (x must be present), t.max\_xor(x) = max over stored y of x \textasciicircum{} y (non-empty). 30 bits.
\begin{lstlisting}
XorTrie t; t.insert(0);  // prefix XORs: best subarray XOR = max over i of t.max_xor(pref[i])
ll best = t.max_xor(x);
\end{lstlisting}
\end{libentry}
\begin{libentry}{KMP (prefix function + pattern search)}{Ctrl+Alt+T, Shift+K}{cpkmp}{above solve()}{O(n + m)}{library/string/kmp.cpp}
prefix\_function(s) -\textgreater{} pi; find\_occurrences(text, pat) -\textgreater{} start index of every match.
\begin{lstlisting}
vector<int> starts = find_occurrences(text, pat);
vector<int> pi = prefix_function(s);
\end{lstlisting}
\end{libentry}
\begin{libentry}{Lyndon factorization (Duval) and minimal rotation}{Ctrl+Alt+U, Y}{cplyndon}{above solve()}{O(n)}{library/string/lyndon.cpp}
duval(s): start indices of the Lyndon factorization s = w1 w2 ... wk with w1 \textgreater{}= w2 \textgreater{}= ... (each wi strictly smaller than all its rotations); min\_rotation(s): start index of the lexicographically smallest rotation.
\begin{lstlisting}
vector<int> starts = duval(s);
int k = min_rotation(s); string canon = s.substr(k) + s.substr(0, k);
\end{lstlisting}
\end{libentry}
\begin{libentry}{Manacher (all palindromes)}{Ctrl+Alt+Y, 6}{cpmanacher}{above solve()}{O(n)}{library/string/manacher.cpp}
manacher(s) -\textgreater{} \{d1, d2\}: d1[i] = odd palindromes centred at i, d2[i] = even ones centred before i; longest\_palindrome(s) -\textgreater{} \{start, len\}.
\begin{lstlisting}
auto [start, len] = longest_palindrome(s);
// number of palindromic substrings = sum(d1) + sum(d2)
\end{lstlisting}
\end{libentry}
\begin{libentry}{Palindromic tree (eertree)}{Ctrl+Alt+U, Shift+E}{cpeertree}{above solve()}{O(n * alphabet) memory, O(n) amortized time}{library/string/eertree.cpp}
Eertree t; t.add(c) for each char (lowercase). Nodes 2.. are the distinct palindromic substrings (len, link = longest proper palindromic suffix); t.last = longest palindromic suffix of the current prefix. occurrences() counts each palindrome\textquotesingle{}s occurrences.
\begin{lstlisting}
Eertree t; for (char c : s) t.add(c);
int distinct = t.t.size() - 2; vector<ll> occ = t.occurrences();
\end{lstlisting}
\end{libentry}
\begin{libentry}{String hashing (mod 2\textasciicircum{}61-1)}{Ctrl+Alt+T, H}{cpstrhash}{above solve()}{O(n) build, O(1) per substring}{library/string/string\_hash.cpp}
StrHash h(s); h.get(l, r) = hash of s[l, r). Random base, hack-safe; same base for every string.
\begin{lstlisting}
StrHash h(s);
if (h.get(l1, r1) == h.get(l2, r2)) { /* s[l1, r1) == s[l2, r2) */ }
\end{lstlisting}
\end{libentry}
\begin{libentry}{Suffix array + LCP}{Ctrl+Alt+Y, 4}{cpsa}{above solve()}{O(n log n)}{library/string/suffix\_array.cpp}
suffix\_array(s): sa[i] = start of the i-th smallest suffix; lcp\_array(s, sa)[i] = LCP of suffixes sa[i] and sa[i + 1].
\begin{lstlisting}
auto sa = suffix_array(s); auto lcp = lcp_array(s, sa);
// distinct substrings = n(n+1)/2 - sum(lcp)
\end{lstlisting}
\end{libentry}
\begin{libentry}{Suffix automaton}{Ctrl+Alt+Y, 5}{cpsam}{above solve()}{O(n * 26)}{library/string/suffix\_automaton.cpp}
SuffixAutomaton sam(s): smallest automaton of all substrings. sam.count\_distinct(), sam.contains(t); states hold len / link / next.
\begin{lstlisting}
SuffixAutomaton sam(s);
ll d = sam.count_distinct(); bool in = sam.contains(t);
\end{lstlisting}
\end{libentry}
\begin{libentry}{Trie}{Ctrl+Alt+T, Shift+T}{cptrie}{above solve()}{O(\textbar{}s\textbar{}) per operation}{library/string/trie.cpp}
Trie t; t.insert(s); t.count(s) = copies of s; t.prefix\_count(p) = words starting with p. a-z only.
\begin{lstlisting}
Trie t;
t.insert(s); t.count(s); t.prefix_count(p);
\end{lstlisting}
\end{libentry}
\begin{libentry}{Z-function}{Ctrl+Alt+T, Z}{cpz}{above solve()}{O(n)}{library/string/z\_function.cpp}
z\_function(s): z[i] = longest common prefix of s and s[i..]; z[0] = 0.
\begin{lstlisting}
vector<int> z = z_function(pat + '#' + text);  // z[i] == |pat| -> match
\end{lstlisting}
\end{libentry}
\section{Geometry}
\begin{libentry}{Closest pair of points}{Ctrl+Alt+Y, 9}{cpclosest}{above solve()}{O(n log n)}{library/geometry/closest\_pair.cpp}
closest\_pair(pts) = smallest squared distance between two of the points (n \textgreater{}= 2), sweep by x with a set by y. Needs cppoint.
\begin{lstlisting}
ll d2 = closest_pair(pts);
\end{lstlisting}
\end{libentry}
\begin{libentry}{Convex hull (monotone chain)}{Ctrl+Alt+T, Shift+V}{cphull}{above solve()}{O(n log n)}{library/geometry/convex\_hull.cpp}
convex\_hull(pts) -\textgreater{} counter-clockwise hull without collinear points. Pt = pair\textless{}ll, ll\textgreater{}.
\begin{lstlisting}
vector<Pt> hull = convex_hull(pts);
\end{lstlisting}
\end{libentry}
\begin{libentry}{Geometry primitives (integer points)}{Ctrl+Alt+T, Shift+3}{cppoint}{above solve()}{O(1)}{library/geometry/point.cpp}
Point \{x, y\} with + - * dot cross norm2; orient(a, b, c) = +1 left turn, -1 right, 0 collinear; on\_segment(p, a, b). Exact.
\begin{lstlisting}
Point a{1, 2}, b, c; cin >> b >> c;
int o = orient(a, b, c); bool on = on_segment(p, a, b); ll d2 = (a - b).norm2();
\end{lstlisting}
\end{libentry}
\begin{libentry}{Half-plane intersection}{Ctrl+Alt+U, 4}{cphalfplane}{above solve()}{O(n log n)}{library/geometry/half\_plane.cpp}
half\_plane\_intersection(hs, B): hs[i] = \{ax, ay, bx, by\} keeps the points LEFT of the directed line a-\textgreater{}b. Returns the intersection polygon (CCW, long double) clipped to the box [-B, B]\textasciicircum{}2, or empty if it is empty.
\begin{lstlisting}
vector<array<long double, 4>> hs = {{0, 0, 1, 0}, {1, 0, 0, 1}, {0, 1, 0, 0}};   // a triangle
auto poly = half_plane_intersection(hs, 1e9);
\end{lstlisting}
\end{libentry}
\begin{libentry}{Minimum enclosing circle (Welzl, randomized)}{Ctrl+Alt+U, Shift+M}{cpmincircle}{above solve()}{O(n) expected}{library/geometry/min\_enclosing\_circle.cpp}
min\_enclosing\_circle(pts): the smallest circle \{cx, cy, r\} containing all points (long double). Random shuffle + incremental construction.
\begin{lstlisting}
auto [cx, cy, r] = min_enclosing_circle(pts);   // pts: vector<pair<long double, long double>>
\end{lstlisting}
\end{libentry}
\begin{libentry}{Polygon: area, point inside, lattice points}{Ctrl+Alt+Y, 8}{cppolygon}{above solve()}{O(n)}{library/geometry/polygon.cpp}
area2(poly) = 2 * signed area (\textgreater{} 0 counter-clockwise); in\_polygon(poly, p): 1 inside, 0 boundary, -1 outside; boundary\_points. Needs cppoint.
\begin{lstlisting}
ll twice = llabs(area2(poly)); int where = in_polygon(poly, p);
ll interior = (twice - boundary_points(poly) + 2) / 2;  // Pick's theorem
\end{lstlisting}
\end{libentry}
\begin{libentry}{Rectangles (intersection, union of two)}{Ctrl+Alt+T, Shift+4}{cprect}{above solve()}{O(1)}{library/geometry/rectangles.cpp}
Rect \{x1, y1, x2, y2\} (x1 \textless{} x2, y1 \textless{} y2): area(), intersect\_area(a, b), union\_area(a, b). USACO bronze staples.
\begin{lstlisting}
Rect a{0, 0, 4, 3}, b{2, 1, 6, 5};
ll both = intersect_area(a, b), either = union_area(a, b);
\end{lstlisting}
\end{libentry}
\begin{libentry}{Rotating calipers (diameter and width of a convex polygon)}{Ctrl+Alt+U, 5}{cpcalipers}{above solve()}{O(n)}{library/geometry/rotating\_calipers.cpp}
For a convex polygon h (CCW, no three collinear, e.g. from the convex hull template): hull\_diameter2(h) = max squared distance between two vertices (exact); hull\_width(h) = min distance between two parallel lines enclosing it.
\begin{lstlisting}
ll d2 = hull_diameter2(hull);   // hull: vector<pair<ll, ll>>
long double w = hull_width(hull);
\end{lstlisting}
\end{libentry}
\begin{libentry}{Segment intersection}{Ctrl+Alt+Y, 7}{cpseginter}{above solve()}{O(1)}{library/geometry/segment\_intersection.cpp}
segments\_intersect(a, b, c, d): closed segments ab and cd share a point (touching and overlapping count). Needs cppoint.
\begin{lstlisting}
bool hit = segments_intersect(a, b, c, d);
\end{lstlisting}
\end{libentry}
\section{Flows \& matching}
\begin{libentry}{Bipartite matching (Hopcroft-Karp)}{Ctrl+Alt+Y, Shift+H}{cphopcroft}{above solve()}{O(E sqrt V)}{library/flow/hopcroft\_karp.cpp}
HopcroftKarp hk(nl, nr); hk.add\_edge(u, v) with u on the left, v on the right; hk.max\_matching(); hk.match\_l[u] = partner or -1.
\begin{lstlisting}
HopcroftKarp hk(nl, nr); hk.add_edge(u, v);
int m = hk.max_matching();  // Konig: min vertex cover = m; max independent set = nl + nr - m
\end{lstlisting}
\end{libentry}
\begin{libentry}{Bipartite minimum vertex cover / maximum independent set (Konig)}{Ctrl+Alt+U, V}{cpkonig}{above solve()}{O(V E) (augmenting paths)}{library/flow/min\_vertex\_cover.cpp}
konig(nl, nr, adj): adj[left] = right neighbours. Returns \{cover\_left, cover\_right\}: a minimum vertex cover, whose size equals the maximum matching; the complement is a maximum independent set.
\begin{lstlisting}
auto [cl, cr] = konig(nl, nr, adj);   // vertices in the cover (bool per side)
\end{lstlisting}
\end{libentry}
\begin{libentry}{Global minimum cut (Stoer-Wagner)}{Ctrl+Alt+U, Shift+A}{cpstoerwagner}{above solve()}{O(n\textasciicircum{}3)}{library/flow/stoer\_wagner.cpp}
stoer\_wagner(w): the minimum total weight of edges whose removal disconnects the graph, w = symmetric n x n weight matrix (n \textgreater{}= 2). No source or sink needed.
\begin{lstlisting}
ll cut = stoer_wagner(w);
\end{lstlisting}
\end{libentry}
\begin{libentry}{Gomory-Hu tree (all-pairs minimum cut)}{Ctrl+Alt+U, A}{cpgomoryhu}{above solve()}{n - 1 max flows (Dinic)}{library/flow/gomory\_hu.cpp}
gomory\_hu(n, edges): a tree par[v], w[v] (v \textgreater{}= 1) on 0..n-1 for an undirected capacitated graph; the min cut between any u and v equals the smallest w on the tree path (Gusfield\textquotesingle{}s algorithm, n - 1 max flows).
\begin{lstlisting}
auto [par, w] = gomory_hu(n, edges);   // edges: {u, v, capacity}
\end{lstlisting}
\end{libentry}
\begin{libentry}{Hungarian algorithm (min-cost assignment)}{Ctrl+Alt+U, Shift+H}{cphungarian}{above solve()}{O(n\textasciicircum{}2 m)}{library/flow/hungarian.cpp}
hungarian(a): a is n x m with n \textless{}= m; assigns each row a distinct column minimizing the total cost. Returns \{cost, col\_of\_row\}. For maximum weight, negate the costs.
\begin{lstlisting}
auto [cost, col] = hungarian(a);   // row i -> column col[i]
\end{lstlisting}
\end{libentry}
\begin{libentry}{Max flow (Dinic) + min cut}{Ctrl+Alt+Y, Y}{cpdinic}{above solve()}{O(V\textasciicircum{}2 E) worst case, O(E sqrt V) on unit-capacity graphs}{library/flow/dinic.cpp}
Dinic g(n); g.add\_edge(u, v, cap); ll f = g.max\_flow(s, t); after it g.min\_cut\_side(v) = v on the source side of a min cut.
\begin{lstlisting}
Dinic g(n); int e = g.add_edge(u, v, cap);
ll f = g.max_flow(s, t); ll flow_on_e = g.flow_on(e);
\end{lstlisting}
\end{libentry}
\begin{libentry}{Max flow with lower bounds}{Ctrl+Alt+Y, Shift+G}{cpflowlb}{above solve()}{one or two Dinic runs}{library/flow/lower\_bounds.cpp}
LowerBoundFlow lb(n); lb.add\_edge(u, v, low, high); lb.max\_flow(s, t) = max s-t flow meeting every bound, -1 if none. Needs cpdinic.
\begin{lstlisting}
LowerBoundFlow lb(n); lb.add_edge(u, v, 1, 3);
ll f = lb.max_flow(s, t);  // -1: the bounds cannot be met
\end{lstlisting}
\end{libentry}
\begin{libentry}{Maximum matching in a general graph (Edmonds\textquotesingle{} blossom)}{Ctrl+Alt+U, Shift+G}{cpblossom}{above solve()}{O(n\textasciicircum{}3)}{library/flow/blossom.cpp}
blossom\_matching(adj): match[v] = partner or -1, for any undirected graph (odd cycles allowed). Matching size = number of v with match[v] != -1, divided by 2.
\begin{lstlisting}
vector<int> mate = blossom_matching(adj);
\end{lstlisting}
\end{libentry}
\begin{libentry}{Min-cost max-flow}{Ctrl+Alt+Y, Shift+Y}{cpmcmf}{above solve()}{O(flow * V * E) worst case, far less in practice}{library/flow/mcmf.cpp}
MCMF g(n); g.add\_edge(u, v, cap, cost); auto [flow, cost] = g.run(s, t, limit). Shortest paths by SPFA, so negative costs are fine.
\begin{lstlisting}
MCMF g(n); g.add_edge(u, v, 1, w);
auto [flow, cost] = g.run(s, t);  // cheapest among maximum flows (or among flows of size limit)
\end{lstlisting}
\end{libentry}
\section{Misc}
\begin{libentry}{De Bruijn sequence}{Ctrl+Alt+U, 9}{cpdebruijn}{above solve()}{O(k\textasciicircum{}n)}{library/misc/de\_bruijn.cpp}
de\_bruijn(k, n): a cyclic sequence of length k\textasciicircum{}n over \{0..k-1\} in which every length-n word appears exactly once (read cyclically). Built from Lyndon words (FKM algorithm). Shortest string containing all words: the sequence + its first n-1 symbols.
\begin{lstlisting}
vector<int> seq = de_bruijn(2, 3);   // 0 0 0 1 0 1 1 1
\end{lstlisting}
\end{libentry}
\begin{libentry}{Interactive problem helpers}{Ctrl+Alt+T, Shift+I}{cpinteractive}{inside solve()}{--}{library/misc/interactive.cpp}
ask(i, j) prints "? i j", flushes, reads the reply (exits on -1); answer(x) prints "! x".
\begin{lstlisting}
int r = ask(i, j);
answer(x);
\end{lstlisting}
\end{libentry}
\begin{libentry}{Random generator}{Ctrl+Alt+T, Shift+R}{cprng}{above solve()}{O(1)}{library/misc/random.cpp}
rnd(l, r) uniform in [l, r]; shuffle(v.begin(), v.end(), rng). Seeded by time.
\begin{lstlisting}
ll x = rnd(1, 100);
shuffle(v.begin(), v.end(), rng);
\end{lstlisting}
\end{libentry}
\begin{libentry}{Random hashing of sets and multisets (XOR / sum hashing)}{Ctrl+Alt+U, X}{cpxorhash}{above solve()}{O(1) per label}{library/misc/xor\_hash.cpp}
RandHash h; h(x) = a random 64-bit label for value x. Multisets are equal (whp) iff the sums of labels match; "every value appears an even number of times" iff the XOR of labels is 0. With prefix sums / XORs this answers range questions in O(1).
\begin{lstlisting}
RandHash h; vector<uint64_t> px(n + 1); for (int i = 0; i < n; i++) px[i + 1] = px[i] ^ h(a[i]);
bool all_even = (px[r] ^ px[l]) == 0;   // range [l, r)
\end{lstlisting}
\end{libentry}
\begin{libentry}{Simulated annealing skeleton}{Ctrl+Alt+U, Z}{cpanneal}{above solve()}{(seconds) x (cost of energy + mutate)}{library/misc/annealing.cpp}
anneal(s, energy, mutate, seconds, t0, t1): minimize energy(state) by random local moves; worse moves are accepted with probability exp(-delta / T), T cooling from t0 to t1 geometrically. Returns \{best state, best energy\}. For heuristic / output-only problems.
\begin{lstlisting}
auto [best, e] = anneal(start, [&](const State& s) { return cost(s); }, [&](State s, mt19937_64& rng) { /* small change */ return s; }, 1.8, 10.0, 0.01);
\end{lstlisting}
\end{libentry}

\end{document}
